\documentclass{article}
\usepackage{graphicx} % Required for inserting images
\usepackage{xcolor} % Must be loaded before soul
\usepackage{soul} % For highlighting text
\usepackage[utf8]{inputenc}
\usepackage{newunicodechar}
\usepackage{tikz}
\usetikzlibrary{positioning}
\usepackage[utf8]{inputenc}
\usepackage[T1]{fontenc}
\usepackage{fontspec}           % For Greek text support with XeLaTeX/LuaLaTeX
\usepackage{polyglossia}        % Multilingual support
\setmainlanguage{english}
\setotherlanguage{greek}

\usepackage{geometry}
\geometry{margin=1in}

\usepackage{csquotes}           % Block quotes
\usepackage{enumitem}           % Better lists
\usetikzlibrary{arrows.meta, positioning, shapes.geometric}

\usepackage{fancyvrb}           % Verbatim environments for ASCII diagrams
\usepackage{setspace}           % Line spacing



% --- Greek support (XeLaTeX) ---
\usepackage{fontspec}
\newfontfamily\greekfont{Gentium Plus} % good polytonic Greek; try Times New Roman if needed
\newcommand{\gk}[1]{{\greekfont #1}}

\newunicodechar{₁}{$_1$}

\usepackage[style=mla]{biblatex} % MLA Citation Style for Bibliography
\addbibresource{references.bib} %Add References page for in text citation
\usepackage[colorlinks=true, linkcolor=blue, urlcolor=blue, citecolor=blue]{hyperref} %enable hyperlinks to be colored and interactable
\usepackage{changepage} %create individual blocks of text that can be altered apart from general document.

\newcommand{\ra}{\ensuremath{\rightarrow}}

\newcommand{\inlinenote}[1]{\par\noindent
  \fcolorbox{orange}{orange!20}{\parbox{0.97\linewidth}{#1}}\par\noindent}

\begin{document}

\section{Emotion: The Form of Desire Under Evaluative Disclosure}


\subsection{E\textit{motion} is \textit{Motion}: The Affective Architecture of Aristotelian Being-Moved (\textit{Paschein})}

When it comes to Aristotle’s account of emotion in the \textit{Rhetoric} (\gk{πάθη}, \textit{pathē}), I am in agreement with Heidegger: ``\gk{πάθη} are not ``psychic experiences,'' are not ``in consciousness,'' (133) nor are they an epiphenomenon sitting atop otherwise rational operations of the human psyche. Rather, as Heidegger continues, they are ``a being-taken of human beings in their full being-in-the-world. That is expressed by the fact that the whole, the full occurrence-context, which is found in this happening, in being-taken, belongs to the \gk{πάθη}. The so-called ``bodily states'' of anxiety, joy, and so forth, are not symptoms, but also belong to the characteristic being of beings, of human beings'' (133). They are a \textit{movement} (kinēsis) of the soul, a dynamic reorientation within the field of appearances that alters how the world shows up and, in so doing, alters how one is disposed to act within it. Emotion is not simply something that happens inside a subject; it is the soul’s \textit{being-moved} by what appears significant to it.\footnote{This anti-internalist, social-ontological reading is the cornerstone of the volume \textit{Heidegger and Rhetoric}; Daniel M. Gross states it with maximal economy: ``Heidegger characterizes \textit{pathos} (variously `passion,' `affect,' `mood,' or `emotion') as the very condition for the possibility of rational discourse, or \textit{logos}. No cynical and crowd-pleasing addition to \textit{logos}, \textit{pathos} is the very substance in which propositional thought finds its objects and its motivation. Without affect our disembodied minds would have no heart, and no legs to stand on'' (Gross, ``Introduction: Being-Moved: The Pathos of Heidegger's Rhetorical Ontology,'' in \textit{Heidegger and Rhetoric}, ed. Gross \& Kemmann, SUNY, 2005, p.~4). Struever's parallel reading anchors the thesis in \textit{kinēsis}: ``Heidegger grasps the emphasis on change, and therefore time. In his account of Aristotelian psychology, \textit{kinesis} gives the authentic \textit{Da-Charakter} to Dasein, the vital place-time of Being'' (Struever, ``\textit{Alltäglichkeit}, Timefulness, in Heidegger's Early Lectures,'' in the same volume, pp.~109--110, citing GA 18, p.~287). Cf. also Withy: ``Every \textit{pathos} that I experience reveals me to myself as moved in some way, positioned and posed by the things I encounter'' (Withy, \textit{Heidegger on Being Affected}, 2023, p.~3).} Emotion is therefore inseparable from \textit{phantasia}—the representational, imagistic, quasi-perceptual capacity through which things “show up” in the first place (DA III.3-8). Indeed, the very word, \textit{emotion}, already encodes the kinetic ontology that Aristotle and Heidegger attribute to \textit{pathos}. According to the Oxford English Dictionary, \textit{emotion} derives from Middle French \textit{émotion} (“a moving, stirring, agitation,” especially in the sense of social or political unrest) and ultimately from Latin \textit{emovēre}, meaning “to move out, remove, agitate.”\footnote{“\textit{Emotion}, \textit{n}.” Oxford English Dictionary, 3rd ed., OUP, 2001. The OED gives the earliest French usage c. 14th century with meanings related to “political agitation or disturbance,” and Latin \textit{emovēre}: “to move out or away; to stir up.”} Morphologically composed of e- (“out, forth”) and movēre (“to move, set in motion”), the term literally designates a moving-out, a displacement or perturbation that carries the subject beyond a prior equilibrium. The modern English form follows this lineage closely, its internal structure—e- (out) + mot- (move) + -ion (denoting the act or process)—naming not a static inner feeling but the \textit{event} of being-moved, the process through which the soul is set into motion in response to what appears significant.\footnote{“\textit{Emotion}, \textit{n}.” OED. See the morphological note under “Etymology”: the noun-forming suffix -ion indicating “action or condition,” attached to \textit{emovēre} via the French \textit{émotion}.} The word’s deeper Indo-European root, *meue- (“to push, move, stir”), likewise links \textit{emotion} to \textit{motion}, \textit{momentum}, and \textit{motivate}, reinforcing that its semantic horizon is fundamentally dynamic.\footnote{“\textit{Move}, \textit{v}.” OED, 3rd ed., OUP, 2002. The entry traces \textit{movēre} to the Proto-Indo-European root *meue-, meaning “to push, move, stir,” which also gives rise to English \textit{motion}, \textit{mobile}, \textit{moment}, etc.}  Thus even at the level of etymology, \textit{emotion} mirrors the Aristotelian claim that \textit{pathos} is a form of \textit{kinēsis}: a being-moved of the soul toward or away from perceived goods and evils.\footnote{Hawhee anchors the \textit{energeia}/\textit{kinēsis} identification in the \textit{Rhetoric} (III.11.1412a.9--10), where Aristotle's reading of Homer culminates in the explicit equation: ``He makes everything move and live, and \textit{energeia} is motion'' --- ``Movement and life, \textit{kinēsis} and \textit{zōos}, are what confer rhetorical presence'' (Hawhee, ``Looking Into Aristotle's Eyes: Toward a Theory of Rhetorical Vision,'' \textit{Advances in the History of Rhetoric} 14.2, 2011, p.~154). The etymological argument I draw from the OED is therefore not merely lexical but tracks an ontological doctrine Aristotle himself underwrites: emotion is the actualized motion of soul-as-life, and the term's morphology preserves what its referent ontologically is.} The \textit{e}- of \textit{emotion} is precisely the ``from without'' of being-taken of Heidegger's formulation, and the ``\textit{motion}'' is the \textit{kinēsis} of being-in-the-world \textit{itself}. To undergo emotion is therefore to be shifted or displaced by the force of appearance already in motion, already reoriented. In this way, the etymology of \textit{emotion} not only confirms but amplifies the ancient insight: emotion is the embodied signature of significance, the soul’s kinetic exposure to what matters. 

However, before I can more fully analyze Aristotelian \textit{pathos} (\gk{πάθος} ``affect''), \textit{pathē} (\gk{πάθη}  ``passions/emotions''), and Heidegger's reading of such, a brief terminological reckoning is required, for Aristotle himself does not deploy \textit{pathos} in a single, univocal sense. The term ranges across at least four distinguishable registers in his corpus. In its broadest, ontological sense, \textit{pathos} is the determination of any being whose constitution admits of alteration---``a quality in respect of which a thing can be altered'' (\textit{Metaphysics} Δ.21, 1022b15)---a sense that ranges over the entirety of natural beings and not merely the affective life of animals. In the aisthetic sense, \textit{pathos} names not just the actualization of the sense organ but the basic hedonic tonality (pleasant/painful) co-given with \textit{aisthēsis}'s own \textit{kritikon}, where ``the soul makes a sort of affirmation or negation, and pursues or avoids the object'' (\textit{De Anima} III.7, 431a8--12). Here I must pause to reiterate the architecture I have previously established: \textit{aisthēsis} and the \textit{orektikon} stand in the ``one substrate, two in being'' relation (\textit{De Anima} III.7, 431a13--14), and the orectic aspect co-given in the perceptual event is \textit{epithymia}---the only species of \textit{orexis} (with \textit{thymos} and \textit{boulēsis}; \textit{De Anima} II.3, 414b1--6) that requires no cognitive enrichment beyond what \textit{aisthēsis}'s hedonic disclosure supplies. Its persisting trace, carried forward by the \textit{resonant kinēsis} after the object has withdrawn, I have coined \textit{resonant epithymia}. In its rhetorical-emotional sense, \textit{pathos} names the determinate higher-order emotions of the \textit{Rhetoric}'s catalog---anger, fear, pity, shame, indignation, envy, emulation, kindness, friendship, and their kin---defined as ``feelings that so change men as to affect their judgements'' (II.1, 1378a21-22). I henceforth call these the \textit{civic passions}: the cataloged species of being-affected (\textit{paschein}) through which the \textit{polis}-dwelling, speech-bearing animal is constituted in its shared world. \textit{Civic passions} marks the contrast with the basic affective valence co-given in \textit{aisthēsis} in that the \textit{civic passions} are the articulationally concretized species that arise from the general affective ground once \textit{doxa}-ratification supplies the propositional content. Lastly, in its bodily-physiological sense, \textit{pathos} names the enmattered alteration through which the formal-hedonic structure of an emotion is realized: ``a boiling of the blood and or warm substance surrounding the heart'' (\textit{De Anima} I.1, 403a30--b1). These senses overlap, shade into one another, and are deployed by Aristotle without consistent demarcation; the Greek is used across the ontological, perceptual, emotional, and physiological registers, and what one says of \textit{pathos} at one register need not, without further argument, hold at another.\footnote{Caston's analysis of \textit{De Anima} I.4, 408b1--18 clarifies what is at stake when Aristotle treats \textit{pathē} as motions: ``Aristotle does not wish to deny that affections in question are motions. On the contrary, he accepts that they are. What he denies is that the soul is the proper subject of these motions \ldots\ the soul is only the terminus and the origin of motions that take place inside the body, not something that itself undergoes these motions'' (Corcilius, ``Aristotle's Definition of Non-Rational Pleasure,'' 2013, p.~54). The four-register distinction drawn here thus tracks an internal articulation of \textit{paschein}-as-motion; it does not commit Aristotle to ``the soul gets angry'' simpliciter (cf. \textit{DA} I.4, 408b13--15: ``it would no doubt be better to say not that the soul pities or learns or thinks, but the man with his soul does so'').}

Before moving forward, I must pause to reiterate that hedonic tonality (good/bad, pleasure/pain, pursue/avoid) is co-given with every moment of being-there: with every perceiving, every thinking, every considering, and, equally, with every \textit{pathos} and \textit{pathē}. Thus, it is to this most basic hedonic tonality to which I refer when I state ``basic affective valence.'' This basic affective valence is one in substrate, two in being with the orectic aspect that arises through it: when the object of sense is actually present, that orectic aspect is \textit{epithymia}; when it persists as part of the \textit{resonant kinēsis} after the object has ceased to be, the residue I have coined is \textit{resonant epithymia}. Perceiving, therefore, elicits \textit{epithymia}, whereas thinking and considering---both of which operate on the \textit{phantasma} of the original object of sense---are accompanied with \textit{resonant epithymia}. \textit{Pathē}, on the other hand, designates the more articulationally concrete emotions cataloged in Aristotle's \textit{Rhetoric}: anger, fear, joy, pity, shame, and their kin. However, both basic affective valence and the \textit{pathē} are treated as two articulational concretions of a single ontological structure: the preservation aspect of \textit{paschein} (``being-affect'') that I discuss below. 

\subsection{\textit{Pathos} and \textit{Pathē}: Two Articulational Concretions}

With our terminological reckoning complete, I can now posit the overall \textit{pathetic} thesis for which I will argue:
\begin{adjustwidth}{0.5in}{0in}
There is a single ontological structure\footnote{Heidegger's ``\gk{ἕξις} (\textit{hexis} ``settled disposition'') is nothing other than a how of \textit{pathos}, being-out-of-composure, in relation to being-composed-as-to'' (BCAP, p. 125) opens the structural space for this two-hows reading. \textit{Pathos} admits of multiple how-modalities; the articulational distinction between basic affective valence and the more articulationally concrete \textit{pathē} is one such how-distinction within the single \textit{pathos} structure.} of \textit{pathos} --- articulated by the four senses of the \textit{Metaphysics} and the three constitutive aspects of the \textit{Rhetoric} --- that is invariant across (a) basic affective valence at the level of perception and \textit{phantasmata}, and (b) the higher-order \textit{pathē} at the \textit{doxa}-articulated level. The difference between the two is not categorical but articulational: (a) realizes the structure at minimal cognitive articulation (the pre-propositional but not \textit{pre-positional} hedonic tonality that accompanies every act of perception); (b) realizes it at maximal cognitive articulation (a more articulationally concrete propositional evaluation under the categories of the \textit{Rhetoric}).\footnote{The present analysis restricts itself to rational animals: those possessing \textit{logos} and therefore capable of \textit{doxa} (``opinion/belief''). Aristotle insists that \textit{doxa} provides the articulational concretion turning the basic affective valence co-given with perception into the higher-order determinacy of shame, pity, and their kin. As Aristotle states: ``there are some of the brutes in which we find imagination, without discourse of reason'' (\textit{De Anima} III.3, 428a23--24). Aristotle elsewhere attributes \textit{pathos}-like states to non-rational animals; the relevant evidence is dense and worth considering at length, since it bears on whether the articulational distinction tracks a real ontological distinction. In the \textit{Nicomachean Ethics}, Aristotle writes: ``Passion also is sometimes reckoned as courage; those who act from passion, like wild beasts rushing at those who have wounded them, are thought to be brave, because brave men also are passionate; for passion above all things is eager to rush on danger\dots Now brave men act for the sake of the noble, but passion aids them; while wild beasts act under the influence of pain; for they attack because they have been wounded or because they are afraid, since if they are in a forest they do not come near one. Thus they are not brave because, driven by pain and passion, they rush on danger without foreseeing any of the perils, since at that rate even asses would be brave when they are hungry; for blows will not drive them from their food; and lust also makes adulterers do many daring things\dots The courage that is due to passion seems to be the most natural, and to be courage if choice and aim be added'' (III.8, 1116b23--1117a9). Here Aristotle simultaneously (a) attributes \textit{pathos}-like states (courage, fear, pain) to non-rational animals and (b) denies them genuine virtue or courage on the ground that they lack the rational element of choice, foresight, and motive: animal ``fear'' and animal ``courage'' are real \textit{pathē}-like dispositions, but they fall short of the higher-order \textit{pathē} because the rational gating is absent. The \textit{Historia Animalium} confirms the pattern from a different direction: ``Of the animals that are comparatively obscure and short-lived the characters are not so obvious to our perception as are those of animals that are longer-lived. These latter animals appear to have a natural capacity corresponding to each of the passions of the soul: to good sense or simplicity, courage or timidity, to good temper or to bad, and to other similar dispositions'' (IX, 588a16--608b). On the present synthesis these attributions are \textit{analogical}: animal ``fear'' is a sustained dispositional orientation grounded in basic valence and persistent \textit{phantasma}, structurally similar to but not identical with the rational \textit{pathē}. The \textit{doxa}-gate is what differentiates the merely-similar from the genuine higher-order \textit{pathē}.}
\end{adjustwidth}
The articulational-concretion distinction I draw between basic affective valence as the perceptual-level minimal articulation of \textit{pathos} and the higher-order \textit{pathē} of the \textit{Rhetoric} as the \textit{doxa}-mediated, maximal articulation fills a gap that Heidegger's treatment of \textit{pathos} in his lectures leaves unaddressed. Heidegger's recovery of \textit{pathos} as a ``being-taken of human beings in their full [bodily] being-in-the-world'' (\textit{BCAP} 133), and the single-structure thesis he extracts from the fourfold sense of the \textit{Metaphysics} and the three constitutive aspects of the \textit{Rhetoric} are analytically and phenomenologically precise. What the undifferentiated treatment cannot accommodate, however, is the variable presence of \textit{pathos} across the action-types Aristotle himself catalogs. The practical syllogism of drinking that Aristotle gives in \textit{De Motu Animalium}---perception, basic appetite, immediate action---does not engage the more articulationally concrete \textit{pathē} Heidegger analyzes, but neither is it \textit{pathos}-less: the perceiver is already disposed toward the perceived object as pursue-able. Without an articulational distinction internal to the \textit{pathos}-structure, this case becomes either a counter-example to the single-structure thesis or a phenomenon that simply falls outside Heidegger's analysis. The articulational distinction resolves both elements: basic affective valence and the higher-order \textit{pathē} are \textit{two articulational concretions of one ontological structure}, the same preservative \gk{σωτηρία} (\textit{sōtēria} ``preservation'')--\textit{paschein} realized at minimal (pre-propositional) and maximal (\textit{doxa}-mediated) articulation.\footnote{The articulational-concretion thesis I am drawing here is consonant with Caston's reading of \textit{phantasia} as the ``common currency among mental states rather than belief,'' which ``emphasizes a form of intentionality which is more basic than the conceptual, and firmly rooted in the general character of perceptual experience'' (Caston, ``Why Aristotle Needs Imagination,'' \textit{Phronesis} 41.1, 1996, p.~52). Caston's pre-conceptual-intentionality reading is consistent with the claim that basic affective valence is already structured by \textit{phantasia}-borne aspectual presentation prior to any \textit{doxa}-mediation; the \textit{doxa}-mediated case is the higher-order articulation of the same underlying intentional structure. Gonzalez's parallel reading in the rhetorical tradition --- that \textit{phantasia} acquires \textit{pistis} only when \textit{doxa} ratifies it --- specifies the articulational threshold at which the higher-order \textit{pathē} of the \textit{Rhetoric} come into determinate being (Gonzalez, ``The Meaning and Function of \textit{Phantasia} in Aristotle's \textit{Rhetoric} III.1,'' \textit{Transactions of the American Philological Association} 136, 2006, pp.~99--131).}

The single-structure thesis without an internal articulational distinction commits one to either of two positions, both of which Aristotle's text resists. On the first position, only the \textit{pathē} of the \textit{Rhetoric} count as genuine \textit{pathē}, and the affective coloring of basic perception is something else: a mere reflex or a physiological tone. The textual evidence appears to resist this: Aristotle insists that ``where there is sensation, there is also pleasure and pain'' (\textit{De Anima} II.2, 413b24), and the structure he names \textit{paschein} extends across the entire range of perceptual and rational activity without internal partition: the same \textit{paschein} that operates in hearing-a-sound operates in being-moved-by-anger, with the relevant difference falling along an internal axis rather than between two adjacent kinds. On the second position, the \textit{pathē} are continuous with basic affective valence in such a way that no real distinction obtains between them, and the determinate emotions of the \textit{Rhetoric} are simply intensifications of the same perceptual coloring. Aristotle's analyses of anger, fear, and pity in the \textit{Rhetoric} do not seem to support this either as they specify propositional contents---the conspicuous slight, the anticipated evil, the undeserved suffering---that the basic affective valence of perception cannot underwrite without \textit{doxa}-mediation. The articulational distinction I am attempting to draw lies in the middle: it preserves the structural unity that the single-structure thesis demands while allowing the differentiated determinacy that Aristotle's per-emotion analyses require. Heidegger's lectures, working at the level of the structural unity, never need to draw the articulational differentiation; the present analysis, working through the chain of actualizing desire, must.\footnote{Frede's recognition that ``Aristotle even provides a subdivision to distinguish between the calculative or deliberative imagination as it functions in human decision-making, and the non-rational imagination that is shared by the animals'' (Frede, ``The Cognitive Role of \textit{Phantasia} in Aristotle,'' in Nussbaum \& Rorty (eds.), \textit{Essays on Aristotle's De Anima}, Clarendon, 1992, p.~287) supplies a parallel structural ground at the \textit{phantasia}-level for the \textit{pathos}-level articulational distinction this section argues for. Where Frede tracks the perceptual/deliberative split in \textit{phantasia} (DA III.10--11), the present analysis tracks the basic-valence/higher-order-\textit{pathē} split in \textit{paschein}: both distinctions are articulational refinements of a single underlying structural register rather than categorical bifurcations into adjacent kinds.}

The unifying ground of this single-structure thesis is the dual \textit{paschein}-distinction Aristotle establishes in \textit{De Anima} has developed canonically as the structural register for perceptual reception: one sense of being-affected is destruction by an opposite, the other is preservation in which a potential is brought into its actuality (II.5, 417b2-7). The present section deploys the second, preservative sense at the level of the higher-order \textit{pathē}. There is no mode of being-there from which pleasure or pain is absent: every moment of our being-there is colored by hedonic tonality. Accordingly, both basic affective valence and the \textit{pathē} are species of \textit{paschein}. The magnitude-axis along which they differ is, however, multi-dimensional. It runs along three independent dimensions: hedonic intensity (how pleasant or painful the affective episode is); existential weight (how deeply the episode bears on the agent's mattering-structure); and cognitive articulation (how propositionally specified the episode's content is). The \textit{pathē} are not simply more intense than basic valence; they are more articulated and---often, though not always---more existentially weighted. Heidegger's reading of the \textit{Metaphysics} fourfold's fourth sense treats magnitude as quantitative-existential, the size of the misfortune that strikes (\textit{BCAP} 132); the conception for which I am arguing preserves this dimension while specifying that the articulational dimension is the one carrying most of the structural weight. The drinking case Aristotle gives in \textit{De Motu Animalium} may involve high hedonic intensity---real thirst is intense---without propositional articulation or substantial existential weight in the rhetorical sense. Conversely, certain reflective episodes may carry high existential weight and propositional articulation without acute hedonic intensity: the considered grief one carries quietly, the determinate shame one acknowledges without active distress. The articulational dimension is, therefore, independent of the intensity dimension: differential articulation is the distinction at issue, not differential strength. What separates basic affective valence from the \textit{pathē} is, accordingly, not a categorical division between two ontological strata but graduation along this multi-dimensional magnitude-axis: it is an \textit{articulational} graduation rather than a difference in kind. Thus, one might walk down a dark alley at night and experience a diffuse uneasiness or anxiety, a bare sense that ``something-should-be-avoided''; and, then, on seeing a hooded figure holding what you perceive to be a weapon (\textit{taking-as}), pass into determinate fear of ``I-am-fearful-of-this-person-now.'' The hedonic tonality is unpleasant in both cases; but, what changes is the episode's articulational concretion: its propositional degree and existential weight. 

\subsubsection{From the Lectures to \textit{Being and Time}}

The articulational distinction drawn here between basic affective valence and the \textit{pathē} is, what I believe to be, the recovery of a structural pattern that Heidegger's own intellectual trajectory presupposes but does not fully articulate at the \textit{pathos} level. Firstly, Heidegger's \textit{Being and Time} retrospectively names Aristotle's \textit{Rhetoric} as the first systematic hermeneutic of attunement (\textit{Stimmung}).\footnote{``The different modes of state-of-mind and the ways in which they are interconnected in their foundations cannot be Interpreted within the problematic of the present investigation. The phenomena have long been well-known ontically under the terms ``affects'' and ``feelings'' and have always been under consideration in philosophy. It is not an accident that the earliest systematic Interpretation of affects that has come down to us is not treated in the framework of ``psychology.'' Aristotle investigates the \gk{πάθη} [affects] in the second book of his \textit{Rhetoric}. Contrary to the traditional orientation, according to which rhetoric is conceived as the kind of thing we ``learn in school,'' this work of Aristotle must be taken as the first systematic hermeneutic of the everydayness of Being-with-one-another. Publicness, as the kind of Being which belongs to the ``they,'' not only has in general its own way of having a mood, but needs moods and ``makes'' them for itself. It is into such a mood and out of such a mood that the orator speaks. He must understand the possibilities of moods in order to rouse them and guide them aright'' (BT \S29, H.138; Eng.~178).} Heidegger here performs an explicit retrospective citation: the \textit{Rhetoric} is the first Aristotelian prefiguration of what \textit{Being and Time} will analyze as attunement (\textit{Stimmung}) and the disclosive structure of disposedness (\textit{Befindlichkeit}). The chronological arc from the lectures \textit{pathos}-analysis to the \textit{Being and Time}'s analysis of disposedness (\textit{Befindlichkeit}) is drawn by Heidegger himself.\footnote{The translation of Aristotelian \textit{pathos} into Heideggerian \textit{Befindlichkeit} has been authoritatively traced by senior Heidegger scholars. Pöggeler states the identification cleanly: ``The speaking person is characterized by \textit{ethos} or attitude (\textit{Haltung}) and by \textit{pathos} or one's feeling of attunement (\textit{Befindlichkeit})'' (Pöggeler, ``Heidegger's Restricted Conception of Rhetoric,'' in \textit{Heidegger and Rhetoric}, 2005, p.~169). Michalski supplies the most explicit formulation of the structural-successor thesis: ``we should think of Aristotle's notion of \textit{pathos}---which Heidegger, let it be noted, translates in his lecture course as `attunement' (\textit{Befindlichkeit}), and which is thereby designated as the historical impulse for Heidegger's own existential-ontological thematization of the attunement of Dasein'' (Michalski, ``Hermeneutic Phenomenology as Philology,'' in the same volume, p.~75). Gadamer's first-person recollection of the 1924 lecture confirms the framing: ``In the 1924 lecture Heidegger especially emphasized `how \textit{logos} had its ground in \textit{pathe}.' The \textit{pathe} are `the fundamental possibilities, by which Dasein primarily orients itself, is attuned\ldots\ The possibility to speak about things is first presented within the so-characterized attunement [\textit{Sich-Befindens}] and being-in-the-world'\ldots'' (Gadamer interview, ``Heidegger as Rhetor,'' in the same volume, p.~57).}

Second, the structural-relation follows the grounding-direction Heidegger establishes in the lectures.
\begin{adjustwidth}{0.5in}{0in}
    The \gk{πάθη} [\textit{pathē}] are not merely an annex of psychical processes, but are rather \textit{the ground out of which speaking arises, and which what is expressed grows back into}; the \gk{πάθη}, for their part, are \textit{the basic possibilities in which being-there itself is primarily oriented toward itself}, finds itself. The primary being-oriented, the illumination of its being-in-the-world is not a \textit{knowing}, but rather a \textit{finding-oneself} that can be determined differently, according to the mode of being-there of a being. (\textit{BCAP} 176).
\end{adjustwidth}
 The companion passage from \textit{Being and Time} makes the structural identity explicit. State-of-mind (\textit{Befindlichkeit}), Heidegger writes, discloses Dasein ``\textit{prior to} all cognition and volition, and \textit{beyond} their range of disclosure'' (BT \S29, H.135; Eng.~175); indeed, ``the mood has already disclosed, in every case, Being-in-the-world as a whole, and makes it possible first of all to direct oneself towards something'' (BT \S29, H.137; Eng.~176). The ``finding-oneself'' of the lectures becomes the disposedness (\textit{Befindlichkeit}) of \textit{Being and Time}. The grounding direction is identical in both texts: the \textit{pathē} ground \textit{logos} in the lectures, and disposedness (\textit{Befindlichkeit}) grounds ``understanding'' and ``discourse'' in \textit{Being and Time}.\footnote{The grounding-direction reading is established in \textit{Being and Time} (\S 29): disposedness (\textit{Befindlichkeit}) is existentially equiprimordial with understanding (\textit{Verstehen}) and discourse (\textit{Rede}), but in the order of disclosure-priority \textit{Befindlichkeit} discloses first: ``a mood assails us. It comes neither from `outside' nor from `inside,' but arises out of Being-in-the-world, as a way of such being'' (BT \S 29, H.136; Eng.~176). The structural finding at the same section: ``in such a state-of-mind\ldots\ Dasein is brought before its Being as `there'\ldots\ this very characteristic of Being which belongs to the `there' is one which we must continue to bring into sharper focus'' (BT \S 29, H.135; Eng.~175). Heidegger's retrospective gloss at \S 29, H.138 (Eng.~178) names Aristotle's \textit{Rhetoric} as ``the first systematic hermeneutic of the everydayness of Being-with-one-another'' --- the precise endorsement of the lecture-text grounding-direction at the level of the \textit{Being and Time} architecture.} Both groundings operate at the most basic level: pre-cognitive, pre-volitional, constitutive of the disclosive medium in which entities can show up at all as mattering. The actualization chain established in the prior section reinstates the same direction at the \textit{pathos}-grade: basic \textit{epithymia} at $A_1$ is pre-doxa, pre-logos, and therefore prior to all cognition and volition; the chain therefore exhibits attunement (\textit{Stimmung}) as the disclosive ground from which \textit{doxa} and the \textit{pathē} proceed.

The fourfold articulation of \textit{pathos} establishes the perceptual case as the structural ground from which the higher-order \textit{pathē} of the \textit{Rhetoric} proceed. Each of the four senses Heidegger gathers from \textit{Metaphysics}---alterability, the \textit{energeia} of alteration, harmful alteration, and the magnitude of the harmful --- was there shown to be satisfied at the perceptual level, with the structural register of preservative \textit{paschein} specifying the manner of its operation. The present section accordingly takes the fourfold as established and turns to its articulationally concrete instantiation at the higher-order \textit{pathē} level: the determinate emotions of anger, fear, and pity, with their propositional contents, conative orientations, and hedonic tonalities. The same single ontological structure specified at minimal articulation in the perceptual case is here realized at maximal articulation (the \textit{doxa}-mediated case); the \textit{pathē} of the \textit{Rhetoric} are not a second \textit{pathos}-structure but a second articulational concretion of the one ontological structure the perceptual case already exhibits.

The four senses, taken as established at the perceptual level, must here be understood as satisfied on both the perceptual and orectic sides: the \textit{resonant aisthēma} is the \textit{energeia} of the \textit{aisthētikon} and \textit{epithymia}/\textit{resonant epithymia} the actuality of the \textit{orektikon}, co-given one-substrate-two-in-being. Heidegger's reading of the fourfold as a ``progressive narrowing'' (\textit{BCAP} 132) specifies the magnitude-axis along which the same structure is realized at varying intensities and articulations: hedonic intensity (how pleasant or painful), existential weight (how mattering), and cognitive articulation (how propositionally specified).

\subsubsection{The \textit{Rhetoric}'s Constitutive Threefold: \textit{Pathos} as Change, \textit{Krisis}, and Hedonic Tonality}

The primary definition of emotion in Aristotle's corpus appears not in \textit{De Anima} or any of the traditionally categorized psychological treatises but in the \textit{Rhetoric}: ``the emotions are all those feelings that so change men as to affect their judgements, and that are also attended by pain or pleasure --- such are anger, pity, fear and the like, with their opposites'' (\textit{Rhetoric} II.1, 1378a20--22). Though the definition appears within the context of rhetorical theory's practical applications, it carries substantial ontological weight; and it is this definition along with the four senses of \textit{pathos} that must anchor any systematic analysis of \textit{pathos} in either of its articulational concretions. In his reading of this same passage from the \textit{Rhetoric}, Heidegger isolates the three constitutive aspects of every \textit{pathos}:
\begin{adjustwidth}{0.5in}{0in}
(1) \gk{Μεταβάλλοντες} [\textit{Metaballontes} ``undergoing change'']: something along the way with respect to which ``a change sets in for us,'' through which ``we change'' from one disposition to another. (2) Combine with this change, \gk{διαφέρουσι πρὸς τὰς κρίσεις}, we ``differentiate ourselves'' from ourselves before the change in that which is the hearer's task: ``to take a position,'' ``to form a view.'' The formation of a view involves the manner and mode in which we change. (3) \gk{οἷς ἕπεται λύπη καὶ ἡδονή}: not ``following,'' but rather ``co-given'' in combination with the \gk{πάθη} is a ``being-disposed-as-higher-or-lower'' of the being-there in question. (\textit{BCAP} 115)
\end{adjustwidth}
\textit{Kriseis}, here, are not external judgments altered by emotion-from-without; they are position-takings of being-there from out of the changed disposition. Heidegger's gloss emphasizes a differentiating-of-oneself-from-oneself, a being-different-from-what-one-was in respect of how one stands toward the matter. Heidegger insists, similarly, that pleasure and pain are not consequences that arrive after the emotion but are co-given with it as a being-disposed-as-higher-or-lower. The \textit{pathos} is the changed-disposition-with-its-position-taking; the pleasure or pain is the tonal elevation or depression co-given with that change. Accordingly, every \textit{pathos}, as articulated in the \textit{Rhetoric} and Heidegger's gloss, is a (1) change-in-disposition that (2) differentiates the subject's position-taking and is (3) co-given with hedonic tonality.

The three-fold articulation is not, however, first introduced at the level of the \textit{pathē}. The second moment in particular---``we differentiate ourselves with respect to position-takings''---is already operative at the perceptual level: Heidegger reads Aristotle's account of \textit{aisthēsis} as a \textit{mesotēs} with the character of \textit{kritikon}, a being-positioned toward possible objects that is itself a capacity for separation and discrimination (\textit{BCAP} 126), with Aristotle himself confirming that the perceiving soul is said to make an ``affirmation or negation'' and to pursue or avoid accordingly (\textit{De Anima} III.7, 431a8--12). This basic affective valence---the organism's pre-propositional orientation toward the pleasant and away from the painful, what Heidegger names the disclosive attunement (\textit{Stimmung}) of disposedness (\textit{Befindlichkeit})\footnote{``What we indicate ontologically by the term `state-of-mind' is ontically the most familiar and everyday sort of thing; our mood, our Being-attuned'' (BT \S29, H.134; Eng.~172).}---is co-given with perception from the earliest moment of the actualization chain. It is dispositional; it colors the perceptual field with the most general affective tonality of good or bad, pursue or avoid. The position-taking moment is therefore constitutive of perception with affective valence; it is not introduced anew by \textit{doxa}. Basic affective valence satisfies the differentiation aspect pre-propositionally but not pre-positionally---as the perceiver's already-being-disposed toward the perceived as pursue-able or avoid-able---whereas \textit{doxa} later re-articulates this position-taking propositionally. The remaining two aspects similarly hold at the basic level. The first---\textit{metaballontes} (``those-undergoing change'')---is satisfied trivially: ``being pained or pleased, or thinking'' are themselves \textit{kinēseis} of the soul (\textit{De Anima} I.4, 408b5--7), and basic affective valence is just this kinetic mode at perceptual actualization. The third---co-givenness of pleasure and pain---is just what basic affective valence is. Emotion (\textit{pathē}), by contrast, is more articulationally concrete. Anger is not undifferentiated displeasure but ``desire accompanied by pain for conspicuous revenge for a conspicuous slight'' (\textit{Rhetoric} II.2, 1378a30). Each emotion, in Aristotle's taxonomy, specifies a propositional evaluation,\footnote{For anger, the slight; for fear, the anticipated evil; for pity, the undeserved suffering (\textit{Rhetoric} II.2, 1378a30--32 for anger; II.5, 1382a21--22 for fear; II.8, 1385b13--15 for pity).} a conative orientation,\footnote{Anger names the conative orientation explicitly: ``desire accompanied by pain for conspicuous revenge'' (\textit{Rhetoric} II.2, 1378a30). For fear the conative orientation is interpretive as Aristotle does not explicitly state it. From the following passage we can surmise---from the formal structure of avoidance---that the conative component is the desire to avoid or flee the anticipated/impending evil: ``Fear may be defined as a pain or disturbance due to imagining some destructive or painful evil in the future. For there are some evils, e.g. wickedness or stupidity, the prospect of which does not frighten us: only such as amount to great pains or losses do. And even these only if they appear not remote but so near as to be imminent: we do not fear things that are a very long way off; for instance, we all know we shall die, but we are not troubled thereby, because death is not close at hand. From this definition it will follow that fear is caused by whatever we feel has great power of destroying us, or of harming us in ways that tend to cause us great pain. Hence the very indications of such things are terrible, making us feel that the terrible thing itself is close at hand; and this---the approach of what is terrible---is danger'' (1382a22-32).} and a hedonic tonality (the pain or pleasure that attends the whole complex). It is this tripartite determinacy that separates emotion from the basic affective valence of pursuit and avoidance. The \textit{Rhetoric}'s three-aspect articulation is therefore not a distinguishing mark of the \textit{pathē} but the general structure of any \textit{pathos}, applicable at every articulational concretion.


\subsection{Enmattered-Accounts}

Aristotle establishes the ontological framework for understanding emotion when he argues that all affections of the soul are ``enmattered accounts'' (\textit{De Anima}, I.1, 403a25-b19) that cannot be abstracted from their material instantiation. The dialectician, Aristotle observes, would define anger as ``a desire for retaliation'' or ``a desire for returning pain for a given pain.'' The physicist, by contrast, would define it as ``a boiling of the blood and hot stuff about the heart'' (403a31-b1). Both are determinate ways of \textit{addressing} anger and bringing it to language. The first yields the \textit{eidos} (the form or ``look'') of anger; the second yields the \textit{hylē} ( the corporeal medium). However, neither account is sufficient without the other; the full definition requires both: anger is the \textit{desire} for retaliation \textit{realized in} the boiling of blood around the heart. In this sense, the \textit{pathē} are meaning-articulations whose meaning is inseparable from their bodily occurrence. The formal and material accounts are complementary rather than competing; indeed, Aristotle argues that anything less would fail to capture the entirety of the phenomenon. Heidegger's remark here is rather telling: ``Only in \textit{phenomenology} has this begun. No division between ``psychic'' and ``bodily acts''!'' (134). Thus, in his description of fear, he states: ``\gk{εἶδος} [\textit{eidos} ``form''] of fearing also has the primary relatedness to a finding-oneself of the body'' (139).\footnote{Gross supplies the most concentrated English rendering of the GA 18 enmattered-\textit{eidos} thesis: ``the \textit{eidos} of fear draws primarily upon a body's condition. The difference lies in the fact that the particular condition of the body (being, say, brown or scratched) plays no role in mathematical inseparability, while for the \textit{pathe} Being in such and such a condition is essential. Both are \textit{logoi enyloi}, but in quite different senses of the term\ldots\ The \textit{eidos} of the \textit{pathe} is a disposition toward other humans, a Being-in-the-world'' (Gross, ``Introduction,'' in \textit{Heidegger and Rhetoric}, 2005, p.~26, translating GA 18, 206--207). The double-articulation (\textit{eidos}-formal-and-bodily) the present paragraph claims is, on Gross's reading, identical with the structural reading Heidegger draws from the dialectician-physicist convergence of \textit{DA} I.1.} Anger is not the desire for retaliation plus, contingently, a hot flush. It is the bodily being-toward-others as having-been-slighted, in which the \textit{eidetic} structure---the dialectician's formal account, comprising both the conative aim at vengeance and the cognitive evaluation of the slight as undeserved---is \textit{already} the form of a bodily occurrence: the eidos and the bodily realization are not two separable contents that the body subsequently expresses, but the single articulated determination of anger's mode of bodily being-toward-others. The shaking and the blood-rush are not symptoms of a hidden judgment; they are part of the way the judgment is \textit{bodily there}. Conversely, the bodily occurrence is not a brute physiological event happening to a soul that judges separately; it is the bodily form of \textit{just this} judgment---the slight, the perpetrator, the wish for retaliation---\textit{articulated} in flesh. It is this double-\textit{articulation} that distinguishes the genuine philosopher of nature who sees the phenomenon as it is: bodily-formed, \textit{eidos}-articulated, world-engaged, being-occurred-to. The \textit{pathē} are not material processes, nor are they immaterial \textit{ideai} or forms; they are structural moments of the way a finite living being finds itself in the world. The implications presented here are considerable. If anger is constitutively enmattered---simultaneously a propositional evaluation (the perceived slight), a conative orientation (the desire to retaliate), a hedonic tonality (the pain from the slight and the pleasure from anticipated vengeance), and a physiological alteration (the heating of blood)---then no analysis or explanation of emotion that reduces it to any one of these components will suffice or contain true explanatory value. Emotion is not merely judgment, though it includes judgment; it is not merely desire, though it includes desire; it is not merely pleasure or pain, though it includes pleasure and pain; it is not merely a bodily change, though it is inseparable from bodily change. The enmattered account thus imposes a structural constraint---which I will explore more thoroughly in a later subsection---on any placement of emotion in the actualization chain: emotion must be located at a point where assertive or propositional, conative, hedonic, and physiological processes converge, and it must be understood as something \textit{irreducible} to any one of them.

\subsection{The Articulational Concretion: From \textit{Doxa} to \textit{Pathē}}

As we have discussed the structural components of which any and all emotions must consist, we must now turn to the conditions required for emotion to be articulationally concretized: ``Further, when we think something to be fearful or threatening, emotion is immediately produced, and so too with what is encouraging; but when we merely imagine we remain as unaffected as persons who are looking at a painting of some dreadful or encouraging scene'' (\textit{De Anima} III.3, 427b21-24). Aristotle's statement that when merely imagining we remain ``unaffected,'' on the surface, seems to be a counter-example: how can the soul be ``unaffected'' if basic affective valence is supposedly always operative? When we look at the painting of a dreadful scene, we are \textit{not entirely} unaffected. We register the painting as more or less compelling, more or less aesthetically pleasing or displeasing, and, thus, basic affective valence is operative. What we lack is the \textit{pathē}: there is no real fear, no real grief. Thus, ``unaffected,'' here, specifies the absence of the more articulationally concrete \textit{pathē}, not absence of all \textit{pathos}.\footnote{Caston's reading of the same passage clarifies what is structurally at stake: ``\textit{phantasia} \ldots\ allows us to conjure up images freely, independent of our views about what is true \ldots\ Because it can function apart from acceptance, moreover, it is possible to remain emotionally unaffected by its contents (3.3, 427a23--24). Belief, in contrast, affects us immediately (427a21--22)'' (Caston, ``Why Aristotle Needs Imagination,'' 1996, pp.~45--46). The unaffected/affected contrast Aristotle draws in \textit{De Anima} (III.3, 427b21--24) is therefore not between absence and presence of \textit{phantasia} but between unratified and ratified \textit{phantasma}: the bare image is \textit{phantasia} without \textit{doxa}; the affecting image is \textit{phantasia} with \textit{doxa}'s assertoric uptake. This is what I call the \textit{doxa}-gate.}\footnote{In \textit{Poetics}, Aristotle explicitly deploys tragedy as a case to argue for \textit{catharsis}: ``A tragedy, then, is the imitation of an action that is serious and also, as having magnitude, complete in itself; in language with pleasurable accessories, each kind brought in separately the parts of the work; in a dramatic, not in a narrative form; with incidents arousing pity and fear, wherewith to accomplish its catharsis of such emotions'' (6, 1449b24-28). If looking at the dramatic representation of dreadful scenes were \textit{entirely apathic}, tragedy could have no cathartic function.} 

While this passage is used by Aristotle to distinguish between \textit{phantasia} and \textit{doxa}, he does so on an affective basis. The bare \textit{phantasma}---the image entertained at will (the fearful scenario merely imagined)---leaves the soul unaffected. This suggests that the more articulationally concretized emotional states (the \textit{pathe} of the \textit{Rhetoric}) derive not from \textit{phantasia} alone. What concretizes them is \textit{doxa}---the involuntary belief that the object is genuinely fearful---and the concretization is immediate, admitting of no perceptible, temporal interval between the formation of the belief/opinion and the onset of emotion. Accordingly, \textit{doxa}, for rational animals, serves as not only a condition but also as a cognitive gate through which emotion must pass to enter the actualization chain. The involuntariness of \textit{doxa} is what gives emotion its distinctive phenomenological character as something that \textit{befalls} the agent rather than something the agent elects to undergo. One does not choose to be afraid; one finds oneself afraid because one has involuntarily \textit{taken-something-to-be} threatening. One does not \textit{choose} to find an object pleasant or painful; we simply \textit{find} them so. Heidegger pinpoints this specifically in his analysis of fear: 
\begin{adjustwidth}{0.5in}{0in}
    The one who is brought into fear mut once ``believe'' that the definite thing that threatens, threatens \textit{him}, and further, that what is threatening proceeds \textit{from this definite human being}, and that he threatens him \textit{now}. What is threatening must be there not only in the sense that I know that one day it could befall me---not being oriented toward the possibility of threat, but rather \textit{being-for, believing,} that I have to expect this or that, that something is now happening to me by this person. (173-174)
\end{adjustwidth}
Without this dispositional uptake---this \textit{believing/opining}-it-to-be-the-case-for-me---fear cannot arise. Heidegger then draws out the negative cases---``those for whom it goes swimmingly'' and ``those who believe that nothing more can happen to them''---``stand outside the possibility of fear'' (174). The lack of fear does not derive from a failure of perception (they may see the threat clearly) but a failure of \textit{appropriating} the threat into \textit{belief about themselves}: ``appropriated as mattering to one'' (174).\footnote{Hyde's reading of \textit{Being and Time} \S\S 54--60 develops a parallel structural claim at the level of conscience: the call of conscience is itself ``a primordial form of \textit{epideictic} rhetoric --- a silent showing-forth (\textit{epideixis}) of the heart of human existence'' (Hyde, ``The Call of Conscience: Heidegger and the Question of Rhetoric,'' in \textit{Heidegger and Rhetoric}, 2005, pp.~94--104). The structural homology between Heidegger's BCAP fear-analysis (the dispositional uptake that makes a threat \textit{appropriated as mattering to one}) and his \textit{Being and Time} conscience-analysis (the silent showing-forth that brings Dasein face-to-face with its ownmost possibility) is what Hyde tracks: in both cases, the affective phenomenon is constituted not by a private inner state but by a structural \textit{taking-on} that mediates between disclosure and uptake. The \textit{doxa}-gate machinery developed here is the Aristotelian articulation of the same structural pattern at the \textit{pathē}-level.}

The individual definitions for various emotions presented in the \textit{Rhetoric} flesh out the position-taking or evaluative structure Aristotle identifies in \textit{De Anima} using anger as his representative example. Each definition, as previously articulated, specifies a cognitive content, a hedonic valence, and---at least implicitly---a conative component. His definition of shame offers a useful explanatory case. ``Shame may be defined as pain or disturbance in regard to bad things, whether present, past, or future, which seem likely to involve us in discredit; and shamelessness as contempt or indifference in regard to these same bad things'' (\textit{Rhetoric} II.6, 1383b15-17). The cognitive evaluation is the judgment of one's own disgrace; that some act, trait, circumstance, etc. (whether already done, happening in the now, or anticipated) will cause others to look upon or judge them unfavorably. The conative component, though not explicitly stated, is easier to parse given the contrastive definition of \textit{shamelessness}. In light of this, the conative element of shame is likely one caring about their standing (``saving face,'' to use a colloquial idiom), and, therefore, the desire to avoid, conceal, or escape whatever threatens them. The shameless person has the same potential cognitive content available (they could recognize the same facts as discrediting) but they lack the motivational care or concern that would make those facts harmful; even though the shameful and shameless person are exposed to the exact same orienting structure the shameful person experiences no dispositional uptake---the belief or opinion that state of affairs is damaging or discrediting---without which shame cannot arise. 

\subsection{Somatic Preparation: Emotion in the Causal Chain}

The causal chain presented by Aristotle in \textit{De Motu Animalium} gives us the clearest illustration of the actualization of desire. To reiterate: sense-perception or thinking \ra imagination \ra desire \ra affections \ra organic parts \ra movement. The location of ``affections''\footnote{The original Greek for ``affections'' here, as translated by Nussbaum, is \gk{πάθη} (\textit{pathē}) (46).} between desire and the organic parts is philosophically significant and one of the leading justifications for my differentiation between basic affective valence---the ``affections'' here being \textit{pathē} in the general, being-affected sense, the somatic substrate (heating, chilling) common to all action---and \textit{pathē} as the more articulationally-concretized, \textit{doxa}-ratified species that arise upon it.\footnote{Nussbaum reads \textit{De Motu Animalium} 8 (702a18 ff.) as establishing the causal chain: ``For the affections suitably prepare the organic parts, desire the affections, and \textit{phantasia} the desire; and \textit{phantasia} comes about either through thought or through sense-perception (702$^{a}$18 ff.). \textit{Phantasia}, then, is involved in every action; it must `prepare' the desire whether or not actual perceiving is going on'' (Nussbaum, ``The Role of \textit{Phantasia} in Aristotle's Explanation of Action,'' in her \textit{Aristotle's De Motu Animalium}, Princeton, 1978/1985, p.~233; cf. p.~221 for the joint-necessity claim of \textit{phantasia} and desire in \textit{De Anima} (432b15 ff.)). The \textit{phantasia}---desire---affections---organic-parts---movement chain is canonical Aristotelian doctrine on this reading; it is consistent with the structural claim that emotion converts \textit{orexis} into determinate-conative form at $M_3 \rightarrow A_4$.}

The causal chain supplied in \textit{De Motu Animalium} suggests that the emotions have a somatic-preparatory role: they are what translate desire into bodily readiness and this is further indicated by the enmattered account and bodily accompaniment. The nature of this readiness is as follows: ``the object we pursue or avoid in the field of action is, as has been explained, the origin of movement, and upon the thought and imagination of this there necessarily follows a heating or chilling. For what is painful we avoid, what is pleasing we pursue, and anything painful or pleasing is generally speaking accompanied by a chilling and heating'' (701b33-702a3). The physiological changes associated with emotion are not incidental accompaniments but the mechanism by which desire is realized in bodily movement. The \textit{affections} at this stage---being-affected in the general, somatic sense---are the preparation of the organic parts; the higher-order emotions are not identical with this stage but are realized \textit{through} it, the \textit{doxa}-ratified species whose somatic component is just such an affection. This convergence between the enmattered account of \textit{De Anima} and the causal chain of the \textit{De Motu Animalium} demonstrates emotion's genuine causal role in the perception-to-movement sequence. The ``boiling of blood around the heart'' that Aristotle identifies as the material component of anger in \textit{De Anima} is one such affection (\textit{pathē} in the general, being-affected sense) that prepares the organic parts for movement in \textit{De Motu Animalium}---the somatic substrate through which the determinate emotion is realized. Furthermore, the fact Aristotle gives the affections a step of their own, distinct from desire, within the causal chain suggests that emotion includes desire without being reducible to it. In each of the emotions discussed, the desire-component is part of the formal structure of the emotion: to define what anger \textit{is} requires mentioning what anger \textit{wants}. Desire (\textit{orexis}) is the moved mover of animal locomotion (\textit{DA} III.10, 433b10-18)---the faculty that, moved by the apparent good, moves the animal. If emotion were simply identical with desire, there would be no reason to distinguish them in the causal chain, and the placement of ``affections'' as a separate stage would be redundant. The ontological status of emotion as separate from desire consists precisely in its composite character: emotion is desire \textit{together with} physiological change \textit{structured by} cognitive content. The desire for retaliation is a component of anger, but anger is the desire for retaliation \textit{as experienced in the boiling of blood, in response to a perceived slight}. It is this composite---not the desire alone---that prepares the organic parts for movement. 

As we have already discussed, the role \textit{doxa} plays in emotion formation reveals that emotion is not a raw feeling or mere physiological disturbance; rather, its is a state of the soul that necessarily involves, as a constitutive condition, a truth-committal cognitive act. The person who fears does not merely \textit{feel} afraid; the person who fears \textit{takes-something-as-truly-fearful-now} fearful, and it is this \textit{taking-as-true} that produces the fear. Indeed, Aristotle's formulation suggests the \textit{doxastic} element is the cognitive component of the emotional state, with the hedonic and physiological components co-given with the truth-commitment or assertion rather than preceding it. However, to reiterate, Aristotle's \textit{doxa} is not a deliberate formation of a considered judgment arrived at through careful reasoning; it is the involuntary taking-of-something-as-being-truly-the-case, a cognitive event that befalls the rational agent rather than one the agent deliberately performs. The speed with which emotion follows \textit{doxa}---``immediately''---suggests that the \textit{doxastic} component of emotion is closer to perceptual recognition than to discursive inference. When the rational agent ``perceiving by sense that the beacon is fire, it recognizes in virtue of the general faculty of sense that it signifies an enemy, because it sees it moving'' (\textit{DA} III.7, 431b5-6), the \textit{doxa} ``this signals danger'' may form almost instantaneously, and fear may follow without any perceptible interval. 

\subsection{The \textit{Pathetic} Loop: Synchronic Chain, Diachronic Saturation}

One highly important aspect of Aristotle's definition of emotion in the \textit{Rhetoric} has yet to be discussed: namely, emotion as ``judgement modifying.'' If a judgment via \textit{doxa} (\textit{taking-as-true}) is necessary for emotion to articulationally concretize, then how can emotion modify judgment? The linear chain $A_0 \rightarrow A_4$ established in prior sections captures the synchronic transitions through which \textit{aisthēsis} grounds the \textit{phantasmata}, which ground the cognitive orientations and \textit{doxa}, \textit{doxa} concretizes the \textit{pathē}, and the \textit{pathē} mobilize action; yet emotion, once concretized, alters the very conditions under which subsequent \textit{doxai} will be formed, under which subsequent judgments are made. Specifically, the chain is not a closed circuit describable through forward transitions alone: each actualization persists as dispositional ground for what follows. This section names that persistence the \textit{pathetic} loop and analyzes its structure on two registers, the synchronic and the diachronic, before grounding the diachronic register in the temporality of \textit{pathos}-grade thrownness.

The seminal text for the feedback's articulation in Aristotle's psychological corpus is \textit{De Insomniis}. Aristotle there observes that: 
\begin{adjustwidth}{0.5in}{0in}
    we are easily deceived respecting the operations of sense-perception when we are excited by emotions, and different persons according to their different emotions; for example, the coward when excited by fear, the amorous person by the object of his desire; so that, with but little resemblance to go upon, the former thinks he sees his foes approaching, the latter, that he sees the object of his desire; and the more deeply one is under the influence of the emotion, the less similarity is required to give rise to these impressions. (460b3-11)
\end{adjustwidth}
The passage isolates a specific cognitive mechanism: under the sway of an active emotion, the controlling sense---the corrective faculty that in ordinary waking cognition intercepts misleading \textit{phantasmata} before they can issue in \textit{doxa}---is impaired, and \textit{phantasmata} that would otherwise be filtered are instead taken as true. Aristotle traces the impairment, indeed, to a structural fact of the cognitive apparatus: the faculty by which the controlling sense judges is not identical with the faculty by which images come before the mind, and the latter can override the former when the perceiving organism is sufficiently agitated. The result is a self-reinforcing recursion. An initial \textit{doxa} concretizes an initial \textit{pathē}; the \textit{pathē} impairs the corrective faculty; subsequent \textit{phantasmata} pass through the corrective gate into \textit{doxa} unopposed; the new \textit{doxai} concretize further \textit{pathē}; and the cycle continues until external intervention, somatic exhaustion, or recovery of the corrective faculty breaks the recursion. The mechanism Aristotle describes is, accordingly, sufficiently general to render intelligible a range of phenomena that are otherwise puzzling. The escalation of anger through sustained rumination on a perceived slight, the intensification of fear through repeated imagining of an anticipated evil, the susceptibility of the dreamer, the ill, and the passionately agitated to \textit{phantasmatic} deception (\textit{De Insomniis} 460b3-16; 461b3-8): each reflects the same loop running unchecked in conditions where the corrective faculty is structurally weakened.\footnote{Caston gives the structural account of why the loop functions this way: ``\textit{Phantasia} is, in effect, an echoing of the initial stimulation in the sense organs: a side-effect, like the original stimulation in character, but unable to compete with fresh, incoming stimulations'' (Caston, ``Why Aristotle Needs Imagination,'' 1996, p.~47, citing \textit{De Insomniis} 2, 459a23--b23). The competition-with-incoming-stimulation structure is what the \textit{De Insomniis} mechanism specifies: when waking-perceptual stimulation is reduced (in dreams) or when emotion lowers the corrective threshold (in agitation, illness, sustained \textit{pathos}), the \textit{phantasma} can override the controlling sense. Corcilius's analysis of animal-motion as a causal-mechanical process develops the parallel reading at the level of bodily \textit{kinēsis}: the same residual-motion structure that grounds \textit{phantasia}'s functional role grounds the somatic feedback through which \textit{pathē} modulate \textit{doxa}-formation (Corcilius, ``Aristotle's Model of Animal Motion,'' 2013, esp.~pp.~52--56).}

Two structural features of the mechanism deserve emphasis. First, the loop is not a malfunction of the cognitive system but an intrinsic feature of the relationship between emotion and \textit{doxa} in rational animals. Aristotle's claim is not that the corrective faculty fails because of a defect to be remedied but that the architecture of judgment is such that emotional intensification induces, of itself, a corresponding lowering of the threshold at which \textit{phantasmata} are \textit{taken-as-true} by \textit{doxa}. Second, the recursion is not symmetric in time: an initial \textit{pathē} lowers the threshold for the \textit{next} \textit{doxa}, which in turn concretizes a \textit{pathē} whose intensity reflects the lowered threshold. Thus the temporal direction of the modulation runs from past to present: prior \textit{pathē} modulate present \textit{doxa}-formation, and present \textit{doxa}-formation feeds forward into present \textit{pathē}. This asymmetry will become decisive when, in the analysis to follow, the feedback is shown to belong to the diachronic rather than the synchronic register of the chain.\footnote{The \textit{pathetic} loop structure identified here is established by Aristotle, as I understand him, across three treatises, each describing the same phenomenon---\textit{pathos} as the structure of being-affected---under a different analytic angle. \textit{De Anima} III.3, 427b21--24 establishes the cognitive condition: \textit{doxa} produces emotion immediately, with no perceptible interval between the truth-committal uptake and the affective response. \textit{De Insomniis} 460b3--16 establishes the recursive mechanism: emotional states permit \textit{phantasmata} to be taken as true when the controlling faculty is impaired, lowering the threshold for further \textit{doxa}-formation. \textit{Rhetoric} II.1, 1378a20--22 establishes the structural consequence: the \textit{pathē} modify \textit{kriseis}, altering the ground on which subsequent judgments take shape. The cross-treatise reading is therefore not a synthesis imposed from outside but the methodological corollary of the claim, defended in \S1.0, that Aristotle's rhetorical ontology must be reconstructed across his treatises rather than read off any one of them---a strategy Heidegger himself follows in the GA 18 lectures.}

Two registers are at work in the \textit{pathetic} loop, and they must be distinguished if the loop's structural function is not to be confused with the chain's directional architecture. The synchronic register concerns the structural transitions at any single moment of \textit{pathos}-actualization: basic affective valence is co-given with perception at $A_1$, carried forward as \textit{resonant epithymia} by the phantasma at $A_2$; \textit{doxa} engages at $A_3$, articulating the \textit{resonant epithymia} into determinate \textit{pathē}; the \textit{pathē} mobilize \textit{orexis} through $M_3 \rightarrow A_4$ into completed action at $A_4$. These transitions are directional and irreversible within a given chain-actualization. The diachronic register, by contrast, concerns how the residues of prior actualizations modulate subsequent actualizations: the residue of a \textit{pathos}-actualization at $T_1$ persists, sedimented in the affected organism's dispositional ground (\textit{hexis}), as a condition shaping the \textit{pathos}-actualization at $T_2$. Accordingly, the \textit{De Insomniis} phenomenon belongs to the diachronic register, not to the synchronic register. The chain remains structurally directional from $A_0$ to $A_4$ within any single actualization; the ``loop'' is the temporal persistence of these residues as ground for subsequent actualizations.

Two distinct mechanisms operate within this diachronic register, and conflating them obscures the precise locus at which prior \textit{pathē} exert their force. The first mechanism---\textit{Stimmung}-saturation---modulates the \textit{input} to \textit{doxa}. A sustained \textit{pathos} persists as a tonal coloring that biases subsequent perception; the cowardly person, indeed, sees enemies in shadows because the fear-\textit{Stimmung} re-saturates the perceptual encounter, so that the basic valence delivered to \textit{doxa} is already shaped by the residue of prior actualizations. The second mechanism---corrective-faculty impairment, the mechanism Aristotle specifies in \textit{De Insomniis}---modulates the \textit{gating} of \textit{doxa}: emotional excitement disables the controlling sense, so that \textit{phantasmata} which would normally be intercepted pass through the corrective gate into \textit{doxa} unopposed. The two mechanisms are commonly conflated under the single rubric of the \textit{pathetic} loop, but they operate at different points of the chain: mechanism (1) at the basic-valence input, mechanism (2) at the corrective gate within the $A_2 \rightarrow A_3$ transition. Both mechanisms are diachronic modulations of existing synchronic transitions, not new transitions added to the chain: on any given pass, the basic affective valence at $A_1$ is already shaped by the residues that prior actualizations have sedimented in the dispositional ground (\textit{hexis}).

This temporal-dispositional persistence is what Aristotle elsewhere analyzes under the doctrine of ἕξις (\textit{hexis}, ``settled disposition'') and what Heidegger names \textit{Geworfenheit}, thrownness; both name a structure in which prior dispositional actualizations persist as the ground from which present disclosure takes its character. The Aristotelian articulation begins with the analysis of habituation in \textit{Nicomachean Ethics}, where Aristotle states that ``moral excellence comes about as a result of habit'' (II.1, 1103a16--17). The claim's significance for the present analysis lies in its specification of \textit{how} the disposition is acquired: through repeated actualization of the chain itself. As Aristotle articulates it, ``for the things we have to learn before we can do, we learn by doing, e.g. men become builders by building and lyre-players by playing the lyre; so too we become just by doing just acts, temperate by doing temperate acts, brave by doing brave acts'' (II.1, 1103a31--b2). Each actualization leaves a residual modification of the dispositional ground; over time, repeated actualizations consolidate into \textit{hexeis} that condition subsequent chain-actualizations. \textit{Hexis} is therefore not a stage of the chain alongside the perceptual, phantastic, cognitive, and orectic transitions; it is the \textit{temporal sedimentation} of repeated chain-actualizations into the dispositional ground that conditions subsequent passes through the chain. The hexis-substrate shapes the resonant-residue content carried by the phantasma at $A_2$, with derived effects of biasing the basic affective valence at $A_1$ (mechanism 1: Stimmung-saturation) and conditioning the gating of \textit{doxa} at the $A_2 \rightarrow A_3$ transition (mechanism 2: corrective-faculty impairment).

The Heideggerian articulation runs structurally parallel. Heidegger specifies in \textit{Being and Time} that disposedness (\textit{Befindlichkeit}) discloses Dasein in three essential respects: it discloses Dasein in its thrownness; it has already disclosed Being-in-the-world as a whole, thereby making it possible to direct oneself toward something at all; and it discloses the worldhood of the world in the mode of circumspective concern (BT \S29, H.137; Eng.~176). Thrownness names, accordingly, the diachronic persistence of prior dispositional states: ``this characteristic of Dasein's Being---this `that it is'---is veiled in its `whence' and `whither', yet disclosed in itself all the more unveiledly; we call it the `thrownness' [\textit{Geworfenheit}] of this entity into its `there'; indeed, it is thrown in such a way that, as Being-in-the-world, it is the `there''' (BT \S29, H.135; Eng.~174). Dasein, in other words, is \textit{always already} disposed; the disposition is not a contingent psychological state superadded to a pre-existing self but the existential-temporal ground from which any present self-relation takes its bearings.

Heidegger's decisive temporal claim, however, is that the \textit{Stimmung} itself temporalizes primarily in the past-modality of \textit{Gewesenheit}: ``Understanding is grounded primarily in the future; one's \textit{state-of-mind}, however, temporalizes itself \textit{primarily} in \textit{having been}. Moods temporalize themselves --- that is, their specific ecstasis belongs to a future and a Present in such a way, indeed, that these equiprimordial ecstases are modified by having-been'' (BT \S68b, H.340; Eng.~390). Thus, the having-already-found-itself-disposed of \textit{Befindlichkeit} is not a past \textit{now stored in memory} but a \textit{Gewesenheit}---a having-been---that \textit{currently} conditions how the entities of the world can come close. The thrown disposition is operative \textit{as} the present disclosure's condition, not as a remembered antecedent of it.

The convergence of the Aristotelian and Heideggerian articulations is, accordingly, no mere analogy: it is the recognition of one phenomenon---prior dispositional actualization persists as the ground from which present disclosure takes its character---described at different ontological registers. Aristotle's articulation operates at the level of perception-and-affection within the soul-structure; Heidegger's operates at the level of Dasein's fundamental disclosedness. The lecture course of 1924 collected as \textit{Basic Concepts of Aristotelian Philosophy} makes the convergence explicit. There Heidegger derives four basic aspects of being-there from the Aristotelian analysis of \textit{hexis}, and the fourth specifies that ``Being-there must, \textit{for itself}, take the opportunity to cultivate this being-composed as a possibility'' (BCAP 122). \textit{Hexis}, Heidegger insists, can have its genesis only ``in being-there itself'' (BCAP 122). The lectures' formulation, indeed, is the prefiguration of what \textit{Being and Time} will name \textit{Geworfenheit}: the dispositional ground that is neither imposed upon Dasein from without nor chosen by it from within but cultivated, in the Aristotelian register, through the very activity that the disposition itself conditions. The \textit{pathetic} loop, accordingly, is not a structurally distinct third direction added to the chain's two synchronic directions; it is the \textit{thrownness-character} of the chain: the way each present actualization is always already shaped by prior actualizations that have sedimented as dispositional ground.

The chain runs synchronically from $A_0$ to $A_4$ within any single moment of \textit{pathos}-actualization, but each moment's $A_1$ is itself the having-been-shaped of prior runs. The actualization chain is, in this precise sense, synchronically directional but diachronically saturated. With this temporal structure now specified---the chain's synchronic directionality and its diachronic temporality both made explicit---the analysis can turn to the question of emotion's specific causal function within the chain: what work, that is, emotion does at the $M_3 \rightarrow A_4$ transition that \textit{orexis} alone could not.

\subsection{Emotion's Function in the Chain}

The coordinated notation for the actualization chain has now been settled: nodes $A_0$ through $A_4$ for the major actualization stages, with $A_3$ as the emotion node (cf. the \textit{pathetic} loop subsection above); and transitions $M_n \rightarrow A_{n+1}$ for the moments between nodes, with $M_3 \rightarrow A_4$ designating the conversion that emotion \textit{is}. What the chain requires at $M_3 \rightarrow A_4$, under evaluatively complex conditions, is the concretization of \textit{orexis} (\gk{ὄρεξις} ``desire'') into a determinate form of desire directed at a specific object under a specific evaluative description. Emotion is precisely this conversion. Without it, the chain would have a generic moved mover---\textit{orexis} set in motion by the apparent good (\textit{De Anima} III.10, 433a31-b1)---but no mechanism to specify \textit{which} form of desire corresponds to \textit{which} disclosure of the apparent good. The desire to flee is not the desire to retaliate; the desire to relieve is not the desire to conceal. Each determinate conative orientation is individuated by the cognitive structure of the emotion that specifies it, and each draws upon a determinate pattern of somatic preparation---the cooling of fear, the heating of anger, the sympathetic pain of pity---that configures the body for the action the emotion calls for.

Emotion does not precede desire as a separate antecedent stage; emotion \textit{is} the form desire takes when \textit{doxa} has disclosed evaluatively complex content. Fear is not something that happens before the desire to flee; fear \textit{is} the desire to flee insofar as it is structured by the cognitive evaluation of an anticipated destructive evil and experienced in the somatic alteration of chilled blood. Anger is not something that happens before the desire to retaliate; anger \textit{is} the desire to retaliate insofar as it is structured by the perception of an unjustified slight and experienced in the somatic alteration of heated blood. Pity is not a separate state that subsequently produces the wish that undeserved suffering cease; pity is the painful recognition of undeserved suffering inseparable from that wish---and similarly for shame, kindness, indignation, and the rest of the \textit{Rhetoric}'s catalog. The \textit{Rhetoric}'s definitions confirm this composite structure: each emotion is a desire of a certain kind, attended by pain or pleasure, directed at a certain kind of object, under a certain evaluative description. The composite is not several parallel components stacked together but a single articulated determination in which the cognitive evaluation, the conative orientation, the hedonic tonality, and the physiological alteration are mutually constitutive aspects of one and the same \textit{pathos}. 

The three-factor schema of \textit{De Anima} III.10 (433a9-b30) is preserved without revision. \textit{Orexis} remains the moved mover, set in motion by the apparent good and itself moving the bodily instrument; the realizable good remains the unmoved originator, drawing the desiring soul without being drawn in turn; the animal via its bodily instrument remains that which is moved, traversing the field of action under the direction of desire. What the schema's original formulation leaves unspecified is the \textit{internal} complexity the moved mover acquires when the apparent good is articulated under the evaluative categories of the \textit{Rhetoric}---when the slighter, the threat, the undeserved sufferer, the public falling-short have entered the soul through \textit{doxa} as evaluatively determinate disclosures of the apparent good. Emotion names this internal complexity. It is the actualization of the \textit{orektikon} (\gk{ὀρεκτικόν} ``faculty of desire'') into a specific form of desire directed at a specific object under a specific evaluative description; it is the form in which the soul's being-moved becomes the body's being-moved; it is the mechanism through which the world's disclosure as mattering in a certain way becomes the animal's motion in the corresponding direction. The schema persists: the moved mover gains an articulational concretion the bare formulation does not require but does not preclude.

The causal function emotion supplies is therefore not universal across the chain but specific to the evaluatively complex case. Where the apparent good is straightforwardly pleasant or painful and the desire is basic appetite or aversion---the drinking case of \textit{De Motu Animalium} (701a32-33), in which appetite, perception or imagination, and immediate action follow without engagement of the \textit{pathē}---basic affective valence alone suffices to direct the organism toward or away from the object. At $M_3 \rightarrow A_4$ in the appetitive case, the chain runs without internal complexity: \textit{orexis} is generic, and the body's being-moved is the immediate consequence of the perception's hedonic valence. The conversion is reserved for the case in which \textit{doxa} has done its articulational work. 

\subsection{Conclusion}

Having worked through both Aristotle's analyses and Heidegger's reading of them, the section's central claim can now be stated in its most concentrated form. Each \textit{pathos}---taking \textit{pathos} as the all-encompassing term for being-affected of the soul, in the broadest sense described in \textit{Metaphysics}---is a determinate concretion of a single five-fold ontological structure. The five-fold, therefore, covers the full register that runs from \textit{epithymia} (basic affective valence in which the object of \textit{orexis} is actively present to the perceiving animal, as water is to the thirsty), through \textit{resonant epithymia} (basic affective valence in which the object is not present and \textit{phantasia} is operating, the \textit{phantasma} substituting for the perceived object)\footnote{It bears reiterating that the orectic side of basic affective valence---\textit{epithymia} when the object is present, \textit{resonant epithymia} when the object is absent and \textit{phantasia} operates on the resonant \textit{kinēsis} to present the phantasma---is co-given throughout with the perceptual side: the \textit{aisthēma} (the perceptual disclosure of the object's form-and-tonality) and \textit{resonant aisthēma} (its persisting trace carried by the phantasma) stand in the ``one substrate, two in being'' relation with the orectic side (\textit{De Anima} III.7, 431a13--14). The two sides cannot be separated; wherever the perceptual disclosure occurs, the orectic aspect is co-given, and vice versa. Each \textit{pathos}-actualization is therefore always already a co-given disclosure of perceptual and orectic content, whether at the minimal articulation of basic affective valence or the maximal articulation of the higher-order \textit{pathē}.}, to the most articulationally concrete \textit{pathē} of the \textit{Rhetoric}---anger, calm, friendship, enmity, fear, confidence, shame, kindness, pity, indignation, envy, emulation---for which this section takes \textit{pathē} and ``emotion'' as interchangeable. The difference across this register are not differences of ontological structure but of \textit{articulational concretion}: the five elements filled maximally in the \textit{pathē} and minimally---but never absently---in basic affective valence. The section's central thesis is, therefore, not a phenomenology proprietary to the \textit{Rhetoric}'s catalog but an ontology of being-affected as such, of which the catalog is the case in which every element of the five-fold is disclosed with the most maximal level of articulation.\footnote{Gross's \textit{Uncomfortable Situations} (Chicago, 2017) develops a parallel social-ontological reading of emotion: ``Aristotle's \textit{Rhetoric} offers alternatives for understanding emotions as social phenomena, which is something that leading humanists like Martha Nussbaum and Richard Sorabji should have remembered as they rushed toward the latest brain science'' (Gross, \textit{Uncomfortable Situations}, 2017, p.~3). And: ``Rhetoric, as a humanistic form of affordance theory, posits a situation that is not neutral but is rather `persuasive,' threatening and promising with respect to our human being-in-the-world'' (p.~20). The five-fold developed here, by integrating the dispositional, intentional, bodily, temporal, and disclosive dimensions, gives the structural-ontological architecture that Gross's transactional-rhetorical reading articulates at the contemporary register; both reject the inner-state model of emotion in favor of a being-in-the-world account.}

The five-fold integrates the three-fold of the \textit{Rhetoric} (dispositional change, \textit{krisis}, hedonic tonality), the four-fold of the \textit{Metaphysics} (alterability, actual alteration, harmful alteration, magnitude), and the four-component enmattered-account analysis of \textit{De Anima} (propositional evaluation, conative orientation, hedonic tonality, physiological alteration), without reducing to any of them. The \textit{Rhetoric}'s \textit{krisis}-moment is recovered by the five-fold as the intentional structure of the second element (the propositional taking-the-matter-as-such-and-such) together with the disclosive function of element 5 (the taking-a-position that this articulation issues in). The hedonic tonality is absorbed into the dispositional moment of element 1, since the way being-there is \textit{taken} in a \textit{pathos} is in part the pleasurable or painful manner of the taking. The \textit{metaphysics}' four-fold gives the underlying ontology of alteration---the bedrock structure in virtue of which \textit{pathos} is the kind of being-affected it is, and from which the \textit{Rhetoric}'s generic ``change'' inherits its analytic content. The enmattered-account contributes most decisively the specification of this ``change'' as \textit{physiological alteration}, which the section develops into element 3, the bodily form: the boiling of blood in anger is not symptomatic accompaniment but constitutive bodily articulation of the \textit{pathos}. What the three integrated schemata together leave underspecified is the temporal articulation of \textit{pathos} (element 4) and the disclosive function through which the \textit{pathos} brings being-there to \textit{logos} (element 5); these are the moments the section has developed and which will now be consolidated. Setting the structure out cleanly, then turning back to its grounding in the disclosive structure of disposedness (\textit{Befindlichkeit}), completes the argument. 

\begin{adjustwidth}{0.5in}{0in}
    \begin{enumerate}
        \item \textbf{The disposition.} Each \textit{pathos} is a determinate manner of \textit{finding-oneself-disposed}. It is not an inner state that the human being inspects; it is the way the human being is taken in this moment, in this situation. Anger is finding-oneself-as-having-been-slighted-by-a-particular-person-who-could-have-acted-otherwise; shame is finding-oneself-as-having-fallen-short-before-an-audience-that-matters; pity is finding-oneself-as-confronted-with-undeserved-suffering-of-one-with-whom-I-could-imagine-myself. The dispositional moment is what the articulational-concretion thesis developed above recovers as the \textit{Befindlichkeit}-side of the \textit{pathos}-structure, and what the \textit{doxa}-gate analysis specifies as the involuntary uptake through which \textit{doxa} concretizes the basic affective valence into determinate emotion; this finding-oneself is somatic-dominant in basic \textit{epithymia} (the tonal coloration of being-thirsty in the present) and \textit{doxastic}-laden in the more articulationally concrete \textit{pathē} (anger's finding-oneself-as-slighted requires the propositional taking-as-true that a slight has occured \textit{and} was undeserved. The difference between the two being one of articulational concretion rather than of underlying structure.
        
        \item \textbf{The toward-which and the toward-whom.} Each \textit{pathos} has a determinate intentional structure: it is \textit{about} this and \textit{with respect to} those. The toward-which articulates the world-side of the disposition (what is at issue); the toward-whom articulates its social-ontological side (the persons to whom and before whom it is articulated). Anger is \textit{with respect to} the slight, \textit{toward} the slighter; pity is \textit{with respect to} undeserved harm, \textit{toward} the one suffering; shame is \textit{with respect to} the fall, \textit{before} the audience that matters. The two together articulate the aspectual disclosure that the \textit{pathos} accomplishes: the way the situation is given as this-and-that, for-me. The per-emotion analyses of the \textit{Rhetoric}'s definitions developed above\footnote{\textit{Rhetoric} II.2, 1378a30--b9 for anger; II.5, 1382a22--32 for fear; II.8, 1385b13--19 for pity; 1383b15--17 for shame.} supplied the specific evaluative contents that fill out this intentional structure for each emotion in the catalog; the \textit{toward-whom} is the element that distinguishes the \textit{pathē} from basic affective valence---thirst has a \textit{toward-which} (the water) but no determinate \textit{toward-whom} in the social-ontological sense, whereas anger, shame, pity, and indignation are constitutively articulated \textit{before} or \textit{toward} determinate others, and the filling of the \textit{toward-whom} slot is what opens the dimension of Being-with (\textit{Mitsein}) on which Heidegger's recovery of the \textit{Rhetoric} turns.\footnote{Heidegger's deliberative-oratorical analysis in \textit{Basic Concepts of Aristotelian Philosophy} identifies the \textit{pathos}-dimension as the structural correlate of the addressee: deliberative oratory turns on ``(1) \textit{ethos}, the `comportment' of the speaker; (2) \textit{pathos}, the `disposition' of the hearers; (3) the matter itself'' (\textit{BCAP}, p.~110). \textit{Pathos}, in other words, is constituted by the affective structure of the hearers' being-disposed-toward-others, not as a private inner state of an isolated subject; it is the structural correlate of the addressee that makes the \textit{Rhetoric}'s catalog a hermeneutic of Being-with rather than a psychology of inner states. The \textit{toward-whom} slot's filling---anger \textit{toward} the slighter, shame \textit{before} the audience that matters, pity \textit{toward} the one suffering---is what fills out this hearer-disposition structure for each emotion in the catalog, opening the \textit{Mitsein}-dimension that the inner-state psychological tradition obscures.}
              
        \item \textbf{The bodily form.} Each \textit{pathos} has a determinate bodily articulation that is constitutive, not symptomatic. The hot rush of anger, the cold of fear, the hot flush of shame, the welling of the eyes in pity --- these are part of the way the \textit{pathos} \textit{is} in the world. To strip a \textit{pathos} of its bodily form is to lose the phenomenon. The bodily form is \textit{logos}-articulated: it has the \textit{eidos} of the slight, the threat, the fall, but it is enmattered, meaning-bodily-bodied-forth. This is the moment that the enmattered-accounts analysis of \textit{De Anima} I.1, 403a25--b19 secured: anger is the desire for retaliation realized in the boiling of blood about the heart, where neither the formal nor the material account is sufficient on its own. The somatic-preparation analysis above specifies the same moment from the action-theoretic side: the bodily form \textit{is} the somatic preparation that translates desire into bodily readiness in the \textit{De Motu Animalium} causal chain (701b33--702a3); the dryness of thirst is as constitutive of basic \textit{pathos} as the boiling blood is of anger, and what varies across the register is not whether the bodily form is constitutive but the degree of \textit{logos}-articulation that overlays the somatic moment.
        
        \item \textbf{The temporal range.} Each \textit{pathos} is structured by a determinate temporality that articulates being-there across all three temporal ecstases. The ecstases are not modes that one \textit{pathos} possesses to the exclusion of others; each \textit{pathos} co-temporalizes \textit{having-been}, \textit{present}, and \textit{projected-futural}, with a primary ecstasis that gives the affect its distinctive phenomenal character. Fear is \textit{primarily} futural --- the coming-toward of an anticipated evil --- but is grounded in the having-been of the dispositional \gk{οἴεσθαι} that lets the threatening matter, and is felt presently as the bodily disquiet of \gk{ταραχή}. Anger is \textit{primarily} having-been --- the slight has already occurred and recoils in present pain --- yet projects forward in the prospective pleasure of vengeance. Shame is, by Aristotle's own explicit formulation, ``pain or disturbance in regard to bad things, \textit{past, present, or future}, which seem likely to involve us in discredit'' (\textit{Rhetoric} II.6, 1383b13--15) --- the three-ecstasis range is named in the definition itself. Pity is \textit{primarily} present --- the confrontation with another's suffering --- but unfolds the past that produced that suffering and the futural projection that the same fate might befall the witness or one of his own. The \textit{pathetic} loop analysis above specifies the temporal structure beneath this characteristic profile: each present \textit{pathos} is always already shaped by the having-been of prior actualizations sedimented as dispositional ground, and Heidegger's \textit{Geworfenheit} names at the existential level what Aristotle's \textit{hexis} names at the perceptual-and-ethical level. The temporal range applies as fully to basic affective valence as to the \textit{pathē}---thirst co-temporalizes having-been (I was thirsty a moment ago, and this conditions my present), present (I am thirsty and see water), and futural projection (drinking will be good, the realizable good being constitutively futural in its very definition) as cleanly as fear braids the three ecstases around its primary futural anchor. 
        
        \item \textbf{The disclosive function: bringing to \textit{krisis}, bringing to \textit{logos}.} Each \textit{pathos} shifts the hearer toward taking-a-position. Each opens or closes specific possibilities of \textit{krisis}, of view-construction or decision. Each cultivates a disposition that, at its limit, brings to speech: anger brings the slighted to confront, fear brings the threatened to confer, shame brings the fallen-short to defend or reform, pity brings the witness to alleviate. The cultivation of the \textit{pathos} is the cultivation of the bringing-to-discourse, and rhetoric, on Heidegger's reading, is the systematic art of just this cultivation. The chain-architectural analysis at $M_3 \rightarrow A_4$ above specifies the mechanism: emotion converts \textit{orexis} into a determinate form of desire under a specific evaluative description, and this determinate desire is precisely what carries the agent into the world-engaging act through which the \textit{pathos} discloses itself in speech and action. Every \textit{pathos} carries a \textit{krisis} (even non-rational animals discriminate \textit{this-as-drink}, but only the \textit{pathē} bring to \textit{logos}---to speech with another---and it is this \textit{logos}-articulation, not the presence of \textit{krisis} as such, that distinguishes them.
    \end{enumerate}
\end{adjustwidth}

These five elements together articulate the generic structure of the \textit{pathē}. Each Aristotelian emotion in the \textit{Rhetoric} is a specific filling-out of the five --- an articulational concretion in which the disposition, the world-side and social-side intentionality, the bodily form, the temporality, and the disclosive function are just so, and not otherwise. The catalog of the \textit{Rhetoric} is therefore not a list of psychological types but a phenomenology of the determinate ways in which a finite, embodied, speaking being finds itself disposed before the matters that concern it. Of the five elements, the temporal range bears the section's heaviest ontological weight and warrants explicit textual development, since the claim that each \textit{pathos} articulates being-there across all three temporal ecstases has direct warrant in both Aristotle and Heidegger and is not a Heideggerian retrojection onto the Aristotelian material.

The Aristotelian textual ground is threefold. First, Aristotle defines memory, perception, and expectation as the three temporal modes of soul's relation to its objects: ``one cannot remember the future --- that is expectation --- or the present --- that is perception; memory relates to what is past'' (\textit{De Memoria} 449b10--15), and ``all memory implies time; therefore only animals that perceive time can remember, and the organ by which they perceive time is that by which they remember'' (449b28--30). Memory and expectation are precisely what \textit{phantasia} makes possible: ``imagination is a feeble sort of sensation, and memory and expectation involve imagination; therefore things remembered and expected can be pleasant'' (\textit{Rhetoric} I.11, 1370a28--33). Because \textit{phantasia} grounds the affective response that follows from \textit{doxa}, and because \textit{phantasia} can range across all three temporal modes, every \textit{pathos} is structurally available to all three ecstases.

Second, Aristotle's per-emotion definitions distribute their primary temporal anchor while preserving traces of all three ecstases. Fear is ``pain or disturbance from an imagination of an evil to come'' (\textit{Rhetoric} II.5, 1382a21): the \gk{μέλλον} names the primary futural ecstasis, yet the very same definition presupposes the present \gk{ταραχή} as bodily disquiet and the having-been of the dispositional \gk{οἴεσθαι} without which, as Aristotle insists, the threat cannot arise into proximity at all (the rich, the powerful, the arrogant ``feel no fear'' precisely because the dispositional ground of believing-oneself-vulnerable has not been laid down). Anger is desire accompanied by pain ``for a conspicuous revenge for a conspicuous slight'' (II.2, 1378a30): the slight has already occurred (having-been-anchor), the pain is felt now (present), and the revenge is anticipated as future pleasure. Indeed, Aristotle remarks, citing Homer, that the anticipation of revenge is ``sweeter by far than the honeycomb'' (1378b1--5, citing \textit{Iliad} XVIII.109). Pity is pain at ``apparent evil befalling one who does not deserve it, and which we might expect to befall ourselves or some friend of ours'' (II.8, 1385b13--15): the witnessing is present, the production of the suffering reaches back to its past causes, and the \textit{phantasmatic} projection ``might befall ourselves'' opens the futural ecstasis through which the witness is implicated in what is witnessed. Third, the shame definition makes the three-ecstasis structure fully explicit: \gk{αἰσχύνη} concerns ``bad things, past, present, or future, which seem likely to involve us in discredit'' (II.6, 1383b13--15) --- Aristotle's own formulation names all three ecstases without remainder.

Heidegger's reading of these passages in GA 18 and his systematic articulation of ecstatic temporality in \textit{Being and Time} provide the ontological framework that recovers what Aristotle has named. In GA 18 \S21, Heidegger glosses fear as ``a disposition that is set before an approaching possibility that pertains to me, comes toward me, and as such announces itself, specifically through the announcement'' (p. 168), and concludes that fear ``discloses the futural character of being-there'' (pp. 168--169). But the future of fear is not a free-floating prospect; it is grounded in the dispositional \gk{οἴεσθαι} --- the believing-oneself-into-danger that being-there has already settled into through prior comportments. This dispositional ground is the \textit{having-been} of fear, and it is what makes courage possible as the \textit{hexis}-correlate that consolidates the past-shaped capacity to fear well: ``Fear has, as a determinate \gk{πάθος}, the possibility of a \gk{ἕξις}. Such a possibility is \textit{courage}. \ldots\ Fear is the condition of the possibility of courage'' (GA 18, p. 175). The threefold structure --- \gk{οἴεσθαι} (having-been-dispositioned), \gk{ταραχή} (present-disquiet), \gk{φοβερόν} as \gk{μέλλον} (futural-threatening) --- is Aristotle's own analytic, recovered by Heidegger as the ecstatic temporal articulation of the affect itself.

\textit{Being and Time} \S68 makes the framework general. Each existential structure has its primary ecstasis: understanding is primarily futural, disposedness (\textit{Befindlichkeit}) primarily grounded in having-been, falling primarily in the present (BT \S68a--c, H.337--348). Yet each ecstasis co-temporalizes with the other two. Fear specifically is analyzed in \S68b as ``a bewildered making-present grounded in forgetting that awaits'' (H.342) --- a temporal structure that braids together making-present (the present-ecstasis of falling), forgetting (the having-been-ecstasis of \textit{Befindlichkeit} in its inauthentic mode), and awaiting (the inauthentic-futural-ecstasis of understanding). The phenomenon of fear, accordingly, is constituted not by a single tense but by the ecstatic unity of all three. The point generalizes to every \textit{pathos}: each is a determinate configuration of how the three ecstases co-temporalize in unity, with a primary ecstasis that gives the affect its distinctive phenomenal character, and the other two ecstases co-active as the ground (having-been) and the projection (futural) that sustain the present feeling-itself of the affect.

The Aristotelian lectures formulate the principle in its most concentrated form: ``Being itself as care and care-full speech is timely; it cares for the still-not-at-hand, speaks about the already-appeared, investigates that which is now with us'' (BCAP, p. 131). What Struever names \textit{Alltäglichkeit} as ``timefulness'' --- the radical specificity of time, the care for tense that defines practical life (Struever, \textit{Heidegger and Rhetoric}, ed. Gross, pp. 110--113) --- is therefore not an external imposition on the \textit{pathē} but their own internal temporal articulation made explicit. The fourth element of the five-fold structure is, accordingly, not a single tense each \textit{pathos} owns but the threefold ecstatic span across which each \textit{pathos} articulates being-there in its determinate way; the primary ecstasis is what gives the affect its name and phenomenal character, while the co-active ecstases supply the conditions of possibility that hold the affect in being.

The five-fold structure has its philosophical center of gravity, accordingly, in the claim that the \textit{pathē} are the ground of \textit{logos}. The point is Heidegger's most concentrated formulation of what the GA 18 lectures recover from Aristotle, and the conclusion's argument turns on it:
\begin{adjustwidth}{0.5in}{0in}
The \gk{πάθη} (\textit{pathē}) are not merely an annex of psychical processes, but are rather \textit{the ground out of which speaking arises, and which what is expressed grows back into}; the \gk{πάθη}, for their part, are \textit{the basic possibilities in which being-there itself is primarily oriented toward itself}, finds itself. The primary being-oriented, the illumination of its being-in-the-world is not a \textit{knowing}, but rather a \textit{finding-oneself} that can be determined differently, according to the mode of being-there of a being. (BCAP, p. 176)
\end{adjustwidth}
Speaking, on this reading, is not a master operation that draws on the \textit{pathē} as resources; speaking arises out of the \textit{pathē} and returns to them, the feedback structure of which the \textit{pathetic} loop analysis above has specified at the \textit{pathos} level itself. The rhetorical art, accordingly, is the cultivation of the dispositional ground that brings being-there to discourse, and Aristotle's \textit{Rhetoric} is, therefore, not a manual of persuasion-techniques but a phenomenology of the disclosive structures through which being-there finds itself in the \textit{polis}.\footnote{Kisiel's reading of GA 18 supplies the matched political-philosophical articulation: ``\textit{Care} is deeply ensconced in the life of \textit{pathos}, mood. The cultivation of the appropriate mood of the auditor by the skilled speaker, so that the attentive listener will then speak with the speaker by speaking after him or her (\textit{Nachreden}) in contagious attunement, suggests that speech finds its deepest roots in mood (GA 18, p.~177). Discourse is not only rational, but impassioned, or better, passional in its rationality \ldots\ As Heidegger puts it, `\textit{ethos} and \textit{pathos} are constitutive of \textit{legein} itself' (GA 18, p.~165)'' (Kisiel, ``Rhetorical Protopolitics in Heidegger and Arendt,'' in \textit{Heidegger and Rhetoric}, 2005, pp.~145--146). The phenomenological-of-the-disclosive-structures reading the present sentence affirms is, on Kisiel's reading, the same structure that grounds Heidegger's political ontology --- and, by extension, the rhetorical protopolitics Aristotle was the first to systematize.} What Heidegger's recovery discloses is that the \textit{pathē} are not psychological phenomena to be added to a prior account of cognition; they are the prior. They are the modes in which being-there has its world disclosed to it, and on the basis of which both cognition and speaking become possible. To analyze anger, fear, shame, pity, indignation, kindness --- the family of higher-order emotions Aristotle catalogs --- is, on this reading, to undertake the phenomenology of disclosedness itself: the inquiry into how a finite, embodied, speaking being finds itself in the world before any explicit knowing or saying, and on the basis of which knowing and saying become possible at all. Heidegger's retrospective citation of the \textit{Rhetoric} as ``the first systematic hermeneutic of the everydayness of Being-with-one-another'' (BT \S29, H.138; Eng.~178) is, in this light, not a passing historical note but a precise diagnostic: the \textit{Rhetoric}'s catalog of \textit{pathē} is the first sustained phenomenology of attunement (\textit{Stimmung}), and the GA 18 lectures are the preparatory work in which that phenomenology is recovered for the \textit{Befindlichkeit}-analytic of \textit{Being and Time}.

The \textit{pathē}, moreover, are co-original with \textit{hexeis}. Both are fundamental concepts of being, articulating the being-structure of beings whose way of being is one in which the \textit{pathē} ``in their way, more proximately determine being-in-the-world, being-in-the-moment'' (BCAP, p.~129) --- they ``characterize the entire human being in its disposition in the world'' (BCAP, p.~129).\footnote{In \textit{Basic Concepts of Aristotelian Philosophy} (\S 18a, p.~129), Heidegger writes that the \textit{pathē} do not concern ``spiritual states'' or ``bodily symptoms'' but characterize the entire human being in its \textit{disposition in the world}; the entire human being is the primary object of the Aristotelian psychology of \textit{De Anima} I, ``a genuine topic not of psychology but of the \textit{discussion of the being of this being}'' (BCAP, p.~129). The co-originality of \textit{hexis} and \textit{pathos} is anchored at the same site at \S 17b-beta (p.~125): ``\textit{hexis} is nothing other than a how of \textit{pathos}, being-out-of-composure, in relation to being-composed-as-to.'' The two together are fundamental concepts of being in the sense that they articulate the being-structure of human existence as practical-rational life lived in speaking-with-one-another.} Each \textit{pathos} has its possible \textit{hexis}-correlate: courage as the \textit{hexis} of fear, good temper (\gk{πρᾳότης} \textit{praotēs}) as the \textit{hexis} of anger, friendliness as the \textit{hexis} of kindness;\footnote{The anger-virtue in the \textit{Nicomachean Ethics} is ``good temper'' (Ross's translation of the \textit{hexis} at IV.5)---the proportionate cultivation of anger rather than its absence. See Sherman, \textit{The Fabric of Character} (Oxford, 1989), and Broadie, \textit{Ethics with Aristotle} (Oxford, 1991), on the \textit{pathos}-\textit{hexis} correlation across the virtues.} \textit{aretē}, in the most general sense, is the cultivation of the \textit{hexis} that fits the \textit{kairos}. The diachronic-register analysis of the \textit{pathetic} loop above specifies the temporal mechanism: \textit{hexeis} are what consolidate over time through the repeated actualization of \textit{pathos}-chains, so that the present \textit{pathos} is always shaped by the dispositional ground that prior actualizations have sedimented. The five-fold structure of any given \textit{pathos} is therefore not a snapshot but a moment within a temporally extended cultivation; the \textit{pathē} are the hinge between the singular event of being-affected and the durable ethical character that consolidates from many such events.

\textit{Hexis}, however, does not operate in a single register across all chain-actualizations; the bivalence requires explicit articulation. Aristotle names \textit{hexis} the \gk{ἐνέργεια} of having (\textit{Metaphysics} Δ.20, 1022b4)---the genuine being-present of having, not a stored possession awaiting deployment---and Heidegger's recovery of this account in \textit{Basic Concepts of Aristotelian Philosophy} \S17 articulates a structural distinction internal to \textit{hexis} itself.

In the case of \textit{technē} (\gk{τέχνη} ``craft''), where the \textit{ergon} (\gk{ἔργον} ``work, product'') is decisive, ``training has the precise sense of \textit{reducing deliberation} insofar as it is through training that the completedness of attaining a result comes about'' (\textit{BCAP}, p.~127). The cultivated \textit{technē-hexis} fires as routine---the prescription ``put further and further out of play'' through repetition---and the trained grammarian writes ``not by chance, but according to a prescription \ldots\ without outside help \ldots\ from out of himself'' (\textit{Nicomachean Ethics} II.4, 1105a22; quoted in \textit{BCAP}, p.~127). The structural analogue at the existential level is the ready-to-hand equipment whose readiness-to-hand consists in its withdrawal from explicit attention: the hammer in the master carpenter's grip disappears into the hammering itself, becoming thematic only when it breaks, goes missing, or stands in the way.\footnote{Heidegger develops the readiness-to-hand analysis in \textit{Being and Time} (\S\S15--16). At \S15: ``The hammering itself uncovers the specific `manipulability' of the hammer. The kind of being which equipment possesses---in which it manifests itself in its own right---we call `readiness-to-hand''' (BT \S15, H.69; Eng.~98). At \S16, equipment becomes \textit{conspicuous} only when it breaks, \textit{obtrusive} when missing, \textit{obstinate} when in the way (BT \S16, H.73--74; Eng.~102--104). The structural homology with \textit{technē-hexis} is precise: trained craft-knowledge operates transparently while the work proceeds, becoming thematic only when the routine fails.} The \textit{technē-hexis}, like the ready-to-hand hammer, withdraws from explicit attention while operating; it becomes thematic only when the routine breaks.

In the case of \textit{praxis} (\gk{πρᾶξις} ``action'') and \textit{aretē} (\gk{ἀρετή} ``virtue''), by contrast, ``the how is only appropriated in such a way that the human being enables himself \textit{to be composed at each moment}; not routine but holding-oneself-open'' (\textit{BCAP}, p.~128). The cultivated \textit{praxis-hexis} sustains rather than reduces the capacity for fresh resolution before the \textit{kairos}; it does not withdraw into transparency but holds the agent open for renewed deliberative engagement on each occasion. The fresh resolution this disposition sustains---the deliberate choice of what to do, here and now, in answer to the \textit{kairos}---is \textit{prohairesis} (\gk{προαίρεσις}, ``deliberate choice''): where the \textit{techn\=e-hexis} contracts deliberation toward routine, the \textit{praxis-hexis} keeps the agent open for renewed prohairetic resolution on each occasion.

These are therefore not species of a single function but structurally distinct modes of being-composed: the former contracts deliberation toward transparent routine; the latter cultivates the standing capacity for deliberative openness. The chain-architectural consequence is direct: where the operative \textit{doxa} is \textit{technē}-grounded, the chain runs through \textit{hexis}-mediated practical inference without engaging the higher-order \textit{pathē}, since the universal premise is craft-procedural; where the operative \textit{doxa} is \textit{praxis}-grounded, the \textit{praxis-hexis} modulates the chain on two registers---biasing the basic affective valence that enters perception, and conditioning which determinate \textit{pathos} the \textit{doxa}-mediated articulational concretion produces.

The disclosive role that the \textit{pathē} discharge in the perception-to-action chain depends on a feature named at several points above but that warrants explicit final statement: emotion responds to appearance, not to corrected judgment. If \textit{phantasia} supplies the appearance, \textit{doxa} is the moment of the soul's commitment to a truth-claim, the unfolding from taking-as-something to taking-something-as-\textit{true}; the \textit{pathos} that follows is a being-moved keyed to the appearance, not to whatever subsequent rational correction the agent may apply. This is precisely what makes the chain run with the speed and inevitability that Aristotle ascribes to it (\textit{De Anima} III.3, 427b21--24). Three illustrative phenomena confirm the analysis. The sound mistaken for a rattlesnake produces real fear, even though the corrected judgment will dissolve the mistake; false beliefs generate real emotions because emotion is keyed to the \textit{phantasmatic} appearance, not to its eventual correction. Art produces emotion, because the \textit{phantasmatic} representation suffices to occasion the affective movement, even when the audience knows the represented event is unreal; and, emotion can contradict subsequent judgment, as when knowing that the rope in the corner is not a snake fails to neutralize the startle, because the rational correction does not unwind the prior \textit{phantasmatic}-doxastic uptake on which the affective movement was built.

Because \textit{phantasia} generates appearances beyond the immediately present, the temporal depth of emotional life follows directly. Aristotle states in \textit{De Memoria} that we remember when ``the movement caused by the object is preserved internally'' (\textit{De Memoria} 450a20--25). This preserved internal movement, when the original object of preservation was an articulationally-concrete emotion---the anger felt at the slight, the joy felt at the pleasure---I coin \textit{resonant pathē}\footnote{\textit{Pathē} in Aristotle's \textit{Rhetoric} II.2--11 catalog (anger, fear, shame, pity, etc.) are articulationally-concrete emotions that I treat throughout as essentially present-tense: a \textit{pathos} is felt in the moment of its actualization, however that actualization is triggered. At \textit{Rhet.} I.11, 1370a25--b18, Aristotle establishes that present-felt pleasure and pain attend both remembering and expecting---the pleasure of anticipated vengeance felt now by the angry man (1370b14), the pleasure of remembered drinks felt now by the fever-thirsty (1370b15--18). Aristotle's ``temporal depth of emotion'' is not, therefore, a distinction between live and non-live emotion but a distinction in what \textit{triggers} the present-felt experience: present perception, future expectation, or memory re-engagement. I coin \textit{resonant pathē} for a structurally distinct phenomenon: the past articulationally-concrete-emotional content preserved alongside the memory itself---the past anger preserved with the memory of the slight, the past joy preserved with the memory of the pleasure. This preserved past-emotional content is the ``movement caused by the object preserved internally'' of \textit{De Memoria} 450a20--25 at the articulationally-concrete-emotional layer. \textit{Resonant pathē} parallels \textit{resonant epithymia} (the past basic-valence-content preserved in the residual \textit{kinēsis}) but at the \textit{doxa}-articulated-emotional layer rather than the basic-affective-valence layer.}: the past \textit{pathos}-content stored in memory itself, distinct from the live \textit{pathos} it can subsequently re-engage. To remember grief is to feel grief again; to anticipate pleasure is to feel pleasure now in relation to a not-yet. When Aristotle writes in the \textit{Rhetoric} that ``both the man who remembers and the man who hopes will be attended by an imagination of what he remembers or hopes'' (\textit{Rhetoric} I.11, 1370a), he articulates the temporal depth of emotional life directly: emotion is not anchored to the present but reaches into past and future by means of the \textit{phantasmatic} appearance. Fear, paradigmatically, is ``a pain or disturbance due to imagining some destructive or painful evil in the future'' (\textit{Rhetoric} II.5, 1382a22--23); the feared event does not exist, the \textit{phantasma} does, and the soul moves as though the feared thing were present. \textit{Phantasia} is therefore not a decorative adjunct to emotion; it is the condition of possibility for the soul's being-moved by absences, possibilities, appearances, and counterfactuals --- which is to say, the condition of possibility for the entire range of higher-order \textit{pathē} placed at the center of evaluatively complex action.

The practical, world-shaping power of emotion that Aristotle's rhetorical theory makes explicit follows from this analysis. To persuade is not primarily to offer proofs but to reshape the listener's \textit{phantasmatic} field. The rhetor must supply images --- concrete, narrative, affectively charged --- that allow the audience to see what the rhetor wants them to see, since this imagistic seeing grounds the subsequent formation of \textit{doxa}, the ensuing movement of emotion, and the eventual act of judgment. If rhetorical persuasion succeeds, accordingly, it is because it sets off the entire cascade of movement: new \textit{phantasma}, new \textit{doxa}, new emotion, new orientation, new action. Heidegger's gloss on this practical-philosophical fact in the lectures is concentrated in the epigraph Gross places at the head of his Introduction: rhetoric is the discipline in which ``the self-elaboration of Dasein is expressly executed'' (Heidegger, SS 1924, GA 18.110, epigraph in Gross ed., \textit{Heidegger and Rhetoric}, p.~1).

For Aristotle, in sum, emotion is the kinetic, affective, evaluative dimension of how the world appears to a finite, embodied, speaking being. It is grounded in \textit{phantasia}, shaped through involuntary \textit{doxa}, felt as pleasure or pain, prepared somatically, and expressed as motion toward or away from what the soul takes to be significant. Emotion is not an add-on to rationality, nor a disruption of cognition, nor a mere physiological disturbance, but the way the soul's self-movement discloses significance at all. The five-fold structure --- disposition, intentionality, bodily form, temporal range, disclosive function --- specifies the architecture of this disclosure for each of the determinate \textit{pathē} the \textit{Rhetoric} catalogs. Emotion, in its full Aristotelian sense, is the affective architecture of being-in-the-world.



\end{document}